% 本文件由 LaTeX 排版 Agent 根据有效 translation.json 自动生成。
% author=Gregory Nazianzen; work_id=3103; title=书信集

\chapter{书信集}

% structure: p0001 source=h1 target=omitted reason=page-title-already-rendered-by-parent

\Emph{第一编} 关于\OriginalTerm{亚波利纳里}{Apollinarian}争论的书信\newline{}\newline{}\Emph{第二编} 与\Person{大圣巴西略}的通信\newline{}\newline{}\Emph{第三编} 杂类书信\par

\SourceRecord{Translated by Charles Gordon Browne and James Edward Swallow.；Nicene and Post-Nicene Fathers, Second Series；Vol. 7.；Edited by Philip Schaff and Henry Wace.；Buffalo, NY: Christian Literature Publishing Co.,；1894.；<http://www.newadvent.org/fathers/3103.htm>.}

% structure: p0001 source=h1 target=section reason=page-title

\section{\WorkTitle{Letters (Division I)}}

\Emph{關於阿波利拿里爭論的書信。}\par

% structure: p0003 source=h2 target=subsection reason=relative-heading-rank; base=h2

\subsection{\Emph{引言}}

這兩封致克肋多紐的書信的背景已在總緒論的第一節中描述過，此處無需多加補充。在致納克塔留（其君士坦丁堡主教位的繼任者）的書信（約公元383年寫成，有時被列為《演講集》第四十六篇）中，聖額我略摘錄了阿波利拿里本人的一部著作，但未提及書名。該論著闡明了此異端的基本謬誤（見緒論第一章，第172頁）。據聖額我略所述，阿波利拿里宣稱天主聖子自永恆就披戴人性身體，而非僅自受孕於童貞瑪利亞之時；但此天主的人性沒有人的理智，其位置由獨生子的神性所代替。他甚至更進一步，將受苦與可死性歸於基督的同一神性。因此，聖額我略鄭重抗議，反對對這些異端者給予任何容忍，甚至不准他們舉行集會；因為他說，容忍或准許必定會被他們視為對其教義立場的寬恕，以及對教會立場的譴責。然而，烏爾曼博士認為，雖然聖額我略無疑說的是真話，即他手中有阿波利拿里的小冊子，但他可能——也許是無意的——誇大了其內容的異端性質，將其論述推向阿波利拿里本會否認的結論。在烏爾曼博士看來，後者的唯一目的是保護基督的合一性；他認為正統的\Quote{兩個完整與完美的性體}的說法傾向於聶斯多略對基督位格的分割；因此他使用了似乎確實混淆了性體、或者至少使降生成人不完美的語言，因為一個缺乏人的理智、其位置由聖子的神性填充的基督，不能說是完美的人。但當厄丕法尼提及此異端的這些極端說法時，懷著對這樣一位好人墮落的惋惜之情（他過去的服務對教會有極大價值），然而，按照當時教會權威常見的精神，他寧願將它們歸因於阿波利拿里的年輕門徒對其教導的擴展，這些門徒並不真正理解阿波利拿里本人的意思。\par

奧林彼烏，本系列最後一封書信的收信人，在公元382年是卡帕多細亞第二省的總督。聖額我略對他極為敬重，並與之保持親密友誼，這從本選集中另一部分額我略致他的其他書信中可以見出。本信的起因是需要向世俗權力求助，以懲罰納齊盎的一群阿波利拿里派信徒，他們竟敢趁聖額我略在克散克薩里溫泉缺席之機，設法祝聖了一位與他們思想一致的主教。技術上講，該主教座堂是空位的，但行政權已由省內各位主教委託給了額我略，儘管他預見到異端者的此類企圖，曾極力敦促府主教及其省內主教們通過合法選舉填補此位，但他絕不準備將自己所受的權威交給一位被非法選舉且不合規矩祝聖的異端者。\par

% structure: p0006 source=h2 target=subsection reason=relative-heading-rank; base=h2

\subsection{\Emph{致君士坦丁堡主教納克塔留。（書信第二百零二封。）}}

天主的眷顾，在过去的岁月中一直守护着教会，如今似乎已完全舍弃了此世生命。我的灵魂因灾祸而深深沉没，以致于我自己生活中的私人痛苦（尽管它们如此众多而巨大，若发生在别人身上，我会认为难以忍受）几乎不值得计入灾难之中；但我只能目睹教会共有的苦难；因为若在当前的危急时刻不努力寻找补救方法，事态将会逐渐陷入完全绝望的境地。那些跟随亚略或欧多克修斯异端的人（我不知是谁煽动他们陷入这种愚妄）正在展示他们的病症，仿佛通过聚集会众获得了一定程度的自信，如同得到许可一般。而马其顿派的人已狂妄到如此地步，竟自称为主教，并在我们的地区游荡，喋喋不休地谈论厄琉息作为他们祝圣的依据。我们内心的祸患——欧诺米，不再满足于仅仅存在；除非他能把所有人都拉入他那毁灭性的异端，否则他自认受了亏待。然而，这一切尚可忍受。教会苦难中最令人痛心的一项，是阿波利拿里派的胆大妄为，不知您为何忽视了，为他们提供了与我平起平坐举行集会的权力。然而，您既由天主的恩宠在神圣奥秘上受到了全面的教导，不仅深知真信仰的辩护，也熟知异端者为反对纯正信仰所发明的一切论证；但或许让尊驾从小人我处听闻也并非不合时宜：有一本阿波利拿里的小册子落入了我手中，其内容超越了所有异端的邪恶。因为他断言，独生子在降生成人为了重塑我们本性时所取的肉身，并非新的获得，而是那肉性的本性从起初就在圣子内。他为这一荒谬断言提出的证据，是从福音中曲解引用的，即\BibleQuote{没有人升过天，除了那从天上降下来的，人子仍旧在天。}\BibleRef{若 3:13}仿佛在他降下之前就是人子，而当他降下时，他带来了那肉身，这肉身似乎在天上就已存在，仿佛从永远就已存在并与他的本质相结合。他又引用宗徒的另一句话，却将其从整个上下文割裂：\BibleQuote{第二个人是出于天的主。}\BibleRef{格前 15:47}于是他假定那从高处降下的人是没有理智的，而独生子的天主性则履行理智的职能，是这人性的第三方部分，因为灵魂和身体属于他的人性方面，但没有理智，理智的位置由天主圣言取代。这还不是最严重之处；最可怕的是他宣称独生子天主——万民的审判者、生命之主、死亡的摧毁者——是有死的，并在其神性本身中受难；在祂身体死亡的三天中，祂的天主性也与身体一同被处死，并由圣父从死者中复活。要逐一列举他在这些荒谬绝伦的主张之外附加的其他命题，将不胜其烦。如今，如果持有如此观点的人有权集会，您在基督内被认可的智慧必须看到：由于我们不赞同他们的观点，任何准许他们集会的许可，无异于宣告他们的观点被认为比我们的更真实。因为若他们被允许像虔敬之人一样教导他们的观点，并满怀信心地宣讲他们的教义，则显然教会的教义已被定罪，仿佛真理在他们一边。因为本性不允许同一主题的两个相反教义都真实。那么，您高贵而崇高的心灵怎能容忍在纠正如此巨大的邪恶时暂时搁置一贯的勇气呢？但即使没有先例，让您无可比拟的完善美德在如此危急时刻挺身而出，教导我们最虔诚的皇帝，若他允许这样的邪恶因言论自由而壮大，以颠覆纯正的信仰，则他在其他方面对教会的热心将毫无益处。\par

% structure: p0008 source=h2 target=subsection reason=relative-heading-rank; base=h2

\subsection{致司铎克莱多尼乌：驳斥阿波利拿里（书信第101封）}

致我们最可敬、蒙天主钟爱的弟兄、同司铎克莱多尼乌，格列高利在主内问安。\par

我想知道，在关于教会的事上，这种革新之风从何而来，竟允许任何人——或如圣经所说，过路的人——撕裂那曾被妥善牧养的羊群，以偷窃式的攻击，或更确切地说，以海盗式和谬误的教导来掠夺它？因为若我们当前的攻击者有任何理由在信仰上谴责我们，即使如此，他们也不应未通知我们就贸然采取此种行动。他们应当先说服我们，或愿意被我们说（若我们至少被视为敬畏天主、为信仰劳苦、帮助教会的人），然后，若有必要，再行革新；但那时或许还有为其暴行辩解的理由。然而，我们的信仰既已在写作中、也在口传中，在此地及远方，在危险与平安之时被宣认，何以有人竟敢如此尝试，而另一些人却保持沉默？\par

其中最为严重之处并不在于（尽管这也令人震惊）这些人通过更坏的人将自己的异端灌输给纯朴的灵魂，而在于他们也对我们说谎，声称我们与他们意见和情感一致；就这样他们挂上饵钩，用这种掩护邪恶地实现自己的意愿，并利用我们的单纯（我们视他们为兄弟而非仇敌）作为他们邪恶的支持。不仅如此，据我所知，他们还宣称已被西方主教会议接纳，而该会议先前曾定他们的罪，这是人所共知的。然而，如果阿波利纳里观点的持有者无论现在或以前曾被接纳，让他们证明此事，我们便满意。因为显而易见，他们只可能是在同意正统信仰的条件下被接纳的，因为除此之外别无可能。他们当然可以证明这一点，要么通过主教会议的记录，要么通过共融函，因为这是主教会议的常规做法。但如果这只是空话，是他们自己的捏造，为了装点门面并借那些人的声望在群众中取得分量，就请教导他们缄默，并驳斥他们；因为我们相信这样的任务非常适合您的生活方式与正统性。不要让他们欺骗自己与他人，声称\Quote{主的人}（他们如此称呼祂，而祂实是我们的主和天主）没有人的理智。因为我们并不将人与天主性分离，而是将\OriginalTerm{位格的合一与同一性}{Unity and Identity of Person}定为教义，祂在古时并非人而是天主，是在万世之前的独生子，不与身体或任何有形之物相混；但在末后的日子里，祂也为了我们的救恩而取了人性；在肉身内可受难，在天主性内不受苦难；在身体上有限，在圣神内无限；同时是尘世的和天上的，可触摸的和不可触摸的，可理解的与不可理解的；为的是通过同一及同一位格——祂既是完美的人也是天主——整个人类因罪而堕落的人性能被再造。\par

若有人不信圣玛利亚是天主之母，他便与天主性隔绝。若有人宣称祂穿过童贞玛利亚如同穿过一个渠道，而不是在她内以神性和人性的方式形成（以神性方式，因为无男人的介入；以人性方式，因为符合孕育的法则），他同样是不虔敬的。若有人主张人性先形成，然后被披上天主性，他也当受谴责。因为这不是天主的诞生，而是逃避诞生。若有人引入两个儿子的观念——一个是圣父的儿子，另一个是母亲的儿子——并否认位格的合一与同一性，愿他失去应许给正确信仰之人的义子名分。因为天主和人乃是两个本性，如同灵魂和肉身是两个；但并非两个儿子或两个天主。因为在此世也没有两个人性；尽管保禄以类似的言语谈到内心与外表的人。我若简洁地说，救主是由彼此不同的元素构成的（因为不可见的与可见的不同，超时间的与受时间限制的不同），然而祂并非两个位格。愿天主不许！因为两个本性因结合而为一，天主性成为人，人性被神化，或随人如何表达。我说不同的元素，因为这与天主圣三中的情形相反；在天主圣三中，我们承认不同的位格，以免混淆位格；但并非不同的元素，因为三位在天主性内是同一的。\par

若有人宣称它如同在先知身上一样，因恩宠运行在祂内，却并非在本体上与祂结合（过去和现在都不是），愿他失去更高能力，或更确切地说，充满相反的能力。若有人不朝拜被钉十字架者，愿他当受诅咒，并被列在杀天主者之中。若有人主张祂因善行而变得完全，或是在祂受洗之后，或是在祂从死者中复活之后，才配得义子名分，如同外邦人所添加的那些加入神祇行列的，愿他当受诅咒。因为那有开始、有进展或被成全的，并非天主，尽管这些说法可用于祂的逐渐显现。若有人主张祂如今已脱下祂的圣肉身，并说祂的天主性失去了身体，且否认祂现今带着身体并与身体一同再来，愿他不得见祂再来的荣耀。因为祂的身体如今在哪里，若不在那取了身体者那里？它并没有被放在太阳下，如同摩尼教徒所妄语的那样，以受侮辱的方式受尊敬；也没有被消散在空气中，如同声音的本性或香气的流动或永不驻留的闪电的轨迹。若是那样，祂复活后如何被触摸，或日后如何被那些刺穿祂的人看见？因为天主性在本质上不可见。不，祂将带着身体再来——我如此学到——如同祂在山上被门徒们看见的那样，或如同祂片刻显现的那样，当祂的天主性压倒了肉体性时。我们说这些是为消除怀疑，而写那些则为纠正新奇的教导。若有人声称祂的肉身是从天降下的，而非来自此地，非出自我们（虽超越我们），愿他当受诅咒。因为经文：\BibleQuote{第二个人出于天，是属天的}（\BibleRef{格前 15:47}），以及\BibleQuote{那属天的怎样，凡属天的也怎样}，以及\BibleQuote{没有人升过天，除了那从天降下而仍在天上的人子}（\BibleRef{若 3:13}），以及类似的话，都应理解为因与天上者的结合而说的；正如\BibleQuote{万有是藉着基督造的}（\BibleRef{若 1:3}）和\BibleQuote{基督住在你们心中}（\BibleRef{弗 3:17}）所说的，不是指着那属于天主的有形本性，而是指着那被心灵察觉的；名称如同本性一样混合，并按照它们亲密结合之律互相融贯。\par

若有人相信祂是一个没有人类理智的人，他实在是没有理智，完全不配得救恩。因为祂所未取的，祂也未治愈；但那与祂的天主性联合的，也得救。若亚当仅仅一半堕落，那么基督所取并拯救的也可能是一半；但若他的全部本性堕落，就必须与那受生者的全部本性联合，如此整个得救。那么，他们不要吝啬给我们完全的救恩，也不要只给救主穿上骨头、神经和人的肖像。因为若祂的人性没有灵魂，连亚略人也承认这一点，以便将祂的受难归于天主性，因为那赋予身体运动的，也是那受苦的。但若祂有灵魂，却无理智，祂如何是人？因为人不是无理智的动物。这必然导致：当祂的形态和帐幕是人性时，祂的灵魂应是马或牛或其他牲畜的灵魂。那么，这就是祂所拯救的；而我便被真理欺骗，被引导夸耀已赐予另一位的尊荣。但若祂的人性是有理智的，并非无理智，那么让他们停止如此真正无理智吧。但有人这样说：天主性取代了人的理智。这与我何干？因为天主性与肉身单独结合不是人，与灵魂单独结合也不是，与两者结合而无理智也不是，理智是人的最主要部分。那么，保存完整的人，并将天主性与之混合，好使你在我整个存在中使我受益。但他断言，祂不能容纳两个完美的性体。除非你只以肉体方式看祂。因为一斗器不能装两斗，一个身体的空位也不能容纳两个或更多身体。但若你注意理智和无形体的事物，请记住：我单一的人格能容纳灵魂、理性、理智和圣神；在我之前，这世界——我指有形无形万物的体系——容纳了圣父、圣子、圣神。因为理智存有的本性就是如此：它们能彼此混合，与形体混合，无形地、不可见地。因为许多声音被一只耳朵容纳；许多眼睛被同一可见对象占据；嗅觉被气味占据；各感官不被彼此限制或挤走，感知物也不因感知之多而减少。但人的理智或天使的理智，有哪一个与天主性相比如此完美，以至于更大者的临在必须挤走另一个？光明与太阳相比算不了什么，一点湿气与河流相比，以至于我们必须先除掉较小的，从房屋拿走光，或从大地拿走湿气，以能容纳更大更完美的。因为一件事物如何能容纳两个圆满？是房屋容纳阳光和太阳，还是大地容纳湿气和河流？这值得探讨；因为这个问题确实值得深思。那么，他们难道不知道，与一物相比是完美的，与另一物相比可能不完美，如同小山与大山相比，或芥菜籽与豆子或其他更大的籽粒相比，尽管它在同类中可称为最大？或者，若你愿意，天使与天主相比，或人与天使相比。所以我们的理智是完美的，有统治权的，但只是相对于灵魂和身体；并非绝对完美；它是天主的仆役和臣属，而非分享祂的元首地位和尊荣。所以梅瑟在法郎面前是天主，\BibleRef{出 7:1} 却是天主的仆人，\BibleRef{户 12:7} 正如经上所载；照耀黑夜的群星被太阳遮蔽，以至于你白天甚至不知其存在；一个小火炬靠近大火，既不被消灭，也不被看见，也不被熄灭；而是成为同一火焰，较大的压倒另一个。\par

但有人可能说：我们的理智受谴责。那么我们的肉身呢？它不也受谴责吗？因此，你或者因罪而弃绝后者，或者因救恩而接纳前者。若祂取了那更差的，好通过祂的降生使之圣化，难道祂不能取那更好的，好通过祂成为人使之圣化吗？如果泥土被发酵成为新面团，哦，智慧的人，难道那肖像不也被发酵，与天主混合，因祂的天主性而被神化吗？我还要加上这一点：若理智被完全弃绝，因其倾向于罪且受定罪，为此祂取了身体却遗漏了理智，那么那些以理智犯罪的人便有托词；因为天主的见证——按你们的说法——表明治愈它是不可能的。让我陈述更重大的结果。你，我的好朋友，侮辱我的理智（你是一个\OriginalTerm{肉体崇拜者}{Sarcolater}，如果我是\OriginalTerm{人崇拜者}{Anthropolater}），好使你将天主束缚于肉体，既然祂无法以别的方式被束缚；因此你拆除了隔墙。但我的理论是什么？我只是一个无知之人，并非哲学家。理智与理智混合，作为更近、更亲近的，并通过它而与肉体混合，成为天主与肉体之间的中保。\par

其次，让我们看看他们对于\Quote{取人性}或如他们所称的\Quote{取肉躯}有何说法。如果这是为了使不可理解的天主成为可理解的，并通过祂的肉躯如同透过帐幔与人交谈，那么他们所扮演的面具和戏剧倒是颇为精巧——更不用说祂本可以如古时在燃烧的荆棘中\BibleRef{出 3:2}以及以人形显现\BibleRef{创 18:5}那样，用其他方式与我们交谈。但如果这是为了通过相似者圣化相似者，以消除定罪，那么正如祂为那已经招致定罪的肉躯而需要肉躯，为我们的灵魂而需要灵魂，同样，祂也需要理智，因为理智不仅在\Person{亚当}内堕落了，而且是首先受影响的，正如医生谈论疾病时所说。因为接受命令者也是未能遵守命令者；未能遵守命令者也是胆敢逾越者；逾越者正是最需要救恩者；而需要救恩者正是祂所取者。因此，理智也被祂取了。这一点现在已通过——用他们自己的说法——几何式的必然证明得到了证实，无论他们是否愿意。但你们的行为就像这样：当一个人的眼睛受伤了，脚也因而受伤时，你们却只顾脚，而任由眼睛无人照管；又或者像一位画家画得不好时，你们却去修改画作，而忽略画家本人，仿佛他成功了。但如果他们被这些论证压倒，便躲进一个命题，即天主即使没有理智也能拯救人，那么我想，祂同样也可以没有肉躯，仅凭意志的举动就能做到，正如祂成就其他一切事，并在没有身体的情况下成就了它们一样。那么，把肉躯连同理智一起拿掉吧，好使你们的荒唐愚蠢臻于完全。但他们被后者蒙蔽了，因此跑到肉躯那里，因为他们不知道圣经的惯用表达。我们也要教导他们这一点。因为凡是知道的人，何须再提起：圣经处处称祂为\Quote{人}和\Quote{人子}？\par

然而，如果他们依据这段经文：\Quote{圣言成了血肉，居住在我们中间}\BibleRef{若 1:14}，并因此抹去人最尊贵的部分（就像鞋匠处理较厚的皮料那样），以便把天主和肉躯连接起来，那么，他们就有理由说天主只是血肉的天主，而不是灵魂的天主，因为经上记载：\Quote{祢赐给祂权柄掌管一切血肉}和\Quote{凡有血肉的都到祢面前来}以及\Quote{愿一切血肉都称颂祂的圣名}——这里所指的就是每一个人。或者，他们又必须设想我们的祖先是没有身体、不可见地下到埃及去的，而且只有\Person{若瑟}的灵魂被\Person{法老}囚禁，因为经上写道：\Quote{他们带着七十五个灵魂下到埃及}\BibleRef{宗 7:14}，以及\Quote{铁器进入了祂的灵魂}——但灵魂是不能被捆绑的。这样辩论的人不明白这些说法使用的是\OriginalTerm{提喻法}{Synecdoche}，即以部分表示整体，正如圣经说雏鸦呼求天主，是指整个鸟类；或者提到昴宿星、金星和北斗星\BibleRef{约 9:9}，是指所有星辰及其眷顾。\par

此外，除非通过提及我们的肉躯，并表明祂为了我们而降至于我们较低的部分，天主对我们的爱就无法以其他方式显明。因为肉躯不如灵魂宝贵，凡稍具理性者都会承认。因此，\Quote{圣言成了血肉}这句话，在我看来，相当于经上说祂\Quote{成为罪}\BibleRef{格后 5:21}或\Quote{成为诅咒}\BibleRef{迦 3:13}为我们——并不是说主变成了这些（怎能如此呢？），而是因为祂承担了它们，从而除去了我们的罪恶，担负了我们的过犯。就目前来说，为了清晰和使众人理解，说这些就够了。我写这些，并非有意撰写一篇论文，而只是要阻止谬误的蔓延；如果认为合适，我将更详尽地阐述这些事。\par

然而，有一件事比这些更为严重，是一个我必须特别提及的要点。我恨不得那搅扰你们的人把自己割掉\BibleRef{迦 5:12}，并且重新引进一个犹太教、一个新的割损礼、一套新的祭献制度。因为如果这样做了，有什么能阻止基督再次降生以废除它们，再次被\Person{犹达斯}出卖、被钉十字架、被埋葬并复活，使一切按相同的次序应验，就像希腊循环周期体系那样，其中相同的星体运转带来相同的事件？因为这些事件中哪些应发生、哪些应省略，其选择方法如何，愿这些以书本众多而自豪的智者向我们说明。\par

但是，由于他们因自己的天主圣三理论而自高自大，错误地指责我们在信仰上不纯正，并引诱群众，有必要让人们知道：\Person{阿波利纳里}虽将\Quote{天主性}之名归于圣神，却未能保存天主性的权能。因为将天主圣三构想为\Quote{大、更大、最大}，如同光、光线和太阳——即圣神、圣子和圣父（这在他著作中明确陈述）——这是一架不指向天堂、而是从天堂降下的天主性阶梯。但我们承认天主圣父、圣子和圣神，这些并非空洞的称号，用以划分等级或权能的不平等，而是正如有同一称号，在天主性中也有同一性体和同一本体。\par

但如果有人以为我们对此主题谈论得正确，却指责我们与异端者共融，就让他证明我们确实应受此指控，我们必会说服他，或退出。但在判决作出之前，作任何革新都是不安全的，尤其是在如此重要且关乎如此重大后果的事上。我们在天主及人面前已经抗议，并继续抗议。而且请放心，即使现在，若不是看见教会正被他们那现今的\OriginalTerm{虚妄会堂}{synagogue of vanity}连同其他诡计撕裂并分裂，我们也不会写这些。但当我们如此说和抗议时，若有人因自己将得的利益、或畏惧人、或极其狭隘的心胸、或牧者及治理者某种疏忽、或喜好新奇与倾向革新，而认为我们不可信，依附于那些人，并分裂教会这高贵身体，那么无论他是谁，他必担受自己的审判\BibleRef{迦 5:10}，并在审判之日向天主交账\BibleRef{玛 12:36}。倘若他们的长篇大论和违反达味圣咏的新\WorkTitle{圣咏集}以及他们韵律的恩宠被视为第三约，那么我们也要作诗篇，并写出大量的韵文。因为我们想我们也拥有天主之神\BibleRef{格前 7:40}，如果这确实是圣神的恩赐，而非人的新奇玩意。我愿你公开宣告这些，免得我们因忽视如此邪恶，并因这邪恶教义从我们的漠不关心获得养分和力量，而负有责任。\par

% structure: p0022 source=h2 target=subsection reason=relative-heading-rank; base=h2

\subsection{驳斥阿波利纳里；致刻勒多尼第二函（第一〇二函）}

既然有许多人前来寻求您的可敬，请求确认他们的信仰，您因此满怀爱心地要求我提出一份关于我见解的简明定义与规则，所以我致函您的可敬，其实您早已知道：我从未且永不可能将任何事物尊崇于尼西亚信经之上，即那次为摧毁亚略异端而聚集的圣教父们的信经；我并且靠天主的助佑，永远持守那信经；详细完备他们关于圣神所说尚不完全之处，因为那问题当时尚未被提出，即我们当相信圣父、圣子及圣神同属一个天主性，如此也明认圣神是天主。请接纳那些如此思考和教导的人共融，如同我也这样行；但那些心持异见的人请拒绝，视他们为天主和天主教会的陌路者。又因为关于人性和降生成人的神圣采纳问题也已提出，请向所有人清楚表明关于我的如下立场：我将由父所生、后来由童贞玛利亚所生的那一位子结合为一，我不称祂为两个子，而是崇拜祂为同一而不可分的天主性和尊荣中的一位。但若有人现在或今后不同意此声明，他将在审判之日向天主交账。\par

现在，我们针对他们关于祂的\OriginalTerm{理智}{Mind}的愚昧见解所反对和坚持的，简而言之如下；因为他们几乎独自坚持他们所设定的条件，仿佛由于缺乏理智，他们才损毁了祂的理智。但为了使他们不指责我们曾经接纳却如今否定他们亲爱的维塔利在可敬的罗马主教达玛苏请求下以书面提交的信仰，我也将对此点作一简短说明。因为这些人在他们真正的门徒和那些被引入他们秘密的人面前进行神学讨论时，如同摩尼教徒在他们称为\Quote{选民}的人当中那样，暴露其疾病的全部程度，几乎不容许救主有任何肉体。但当他们被关于降生成人的圣经所呈现的普遍回答所驳斥和追问时，他们确实承认正统言辞，却曲解其含义；因为他们承认那人性既非无灵魂、亦非无理性、亦非无理智、亦非不完全，但他们却引入天主性来取代灵魂、理性和理智，好像天主性仅仅与祂的肉身混合，而未与我们人类的其他属性混合；尽管祂的无罪远超越我们，且是我们情欲的净化。\par

因此，他们错误地解释\BibleQuote{我们却有基督的理智}这句话\BibleRef{格前 2:16}，而且极其荒谬地，当他们说祂的天主性就是基督的理智，并不像我们那样理解那段落，即那些通过模仿救主从我们取来的理智而净化了自己理智的人，并已可能达到与它一致的人，就被说成拥有基督的理智；正如那些训练了自己的肉体，并在这方面成为同一身体和有分于基督的人，可被证为拥有基督的肉体；所以他说：\BibleQuote{我们怎样带了土的肖像，也要怎样带天上的肖像}\BibleRef{格前 15:49}。因此他们宣称\Quote{完全的人}并非那位在各方面都受过试探、像我们一样，却无罪的\BibleRef{希 4:15}，而是天主与肉身的混合。因为，他们说，还有什么比这更完全的呢？\par

他们对待描述降生成人的词也使用同样的诡计，即：\Quote{祂成了人}，却将其解释为并非指祂以所取的人性成为人——按照圣经，\BibleQuote{祂知道人心里所存的}\BibleRef{若 2:25}——而是教导说，这意味着祂与人类交往和谈话，并托庇于那句经文，说\Quote{祂在地上被看见，与人交往}\BibleRef{巴 3:37}。还有人能再争辩什么呢？那些既剥夺人性又剥夺内在肖像的人，用他们新发明的面具只清洁我们的外表\BibleRef{玛 23:25}以及可见的部分；他们如此自相矛盾，以致于有时候为了肉身的缘故，以粗俗和属肉体的方式解释其余一切（因为他们由此衍生出他们的\OriginalTerm{第二犹太主义}{second Judaism}、在乐园里千年的愚蠢享乐，以及几乎认为我们在这千年之后将再次恢复同样状态的想法）；而另一些时候他们又把祂的肉身说成是\OriginalTerm{幻影}{phantasm}而非实体，仿佛祂没有经历任何我们的体验，甚至连那些无罪者所经历的也没有；为此他们使用宗徒的话，以一种非宗徒的方式理解和陈述，说\Quote{我们的救主取了人的样式，也取了人的形状}\BibleRef{斐 2:7}，仿佛这些话所表达的不是人的形状，而是某种欺骗性的幻影和外观。\par

既然这些说法，若正确理解则支持正统，若错误解释则是异端，那么如果我们以更正统的意义接受\Person{Vitalius}的话语，有什么可惊讶的呢？我们渴望它们如此被理解，这说服了我们，尽管其他人对他著作的意图感到愤怒。我认为，这就是\Person{Damasus}本人后来得到更好信息，同时得知他们坚持先前的解释，因此将他们逐出教会，并以\OriginalTerm{绝罚}{Anathema}推翻他们书面的信仰宣认的原因；同样也是由于他对因单纯而遭受的欺骗感到恼怒。\par

既然他们已公开被证实此事，但愿他们不要发怒，而应自愧；不要诽谤我们，而要谦卑自己，并从他们的门廊擦去那伟大而惊人的宣告和对自己正统的夸耀，以那问题和区分迎接所有进入的人：我们应敬拜的，不是一位\OriginalTerm{承载天主的人}{God-bearing Man}，而是一位\OriginalTerm{承载肉身的天主}{flesh-bearing God}。还有什么比这更不合理的呢？尽管这些真理的新宣扬者很看重这个称号。因为虽然它因其对立的敏捷而具有某种诡辩的优雅，以及一种对未受教育者来说讨喜的戏法般的夸耀，但它却是荒谬中最荒谬的，愚蠢中最愚蠢的。因为如果有人将\Quote{人}或\Quote{肉身}改为\Quote{天主}（前者我们喜欢，后者他们喜欢），然后使用这奇妙的、被神圣认可的对立，我们会得出什么结论呢？即我们应敬拜的，不是一位\OriginalTerm{承载天主的肉身}{God-bearing Flesh}，而是一位\OriginalTerm{承载人的天主}{Man-bearing God}。啊，怪异的荒谬！他们今天向我们宣告一种从基督时代以来一直隐藏的智慧——一件值得我们流泪的事。因为如果信仰始于三十年前，而此时自从基督显现以来已近四百年过去了，那么我们的福音在这期间就是枉然的，我们的信仰也是枉然的；殉道者将白白作证，如此众多且伟大的主教们将白白治理百姓；而恩宠成了韵律的事，而非信德的事。\par

谁会不惊叹他们的学识呢？因为他们凭自己的权威分割基督的事，将祂出生、受试探、饥饿、口渴、疲倦、睡眠之类的话归于祂的人性；而将祂被天使荣耀、战胜试探者、在旷野喂养百姓、以那样的方式喂养他们、以及在水面上行走之类的话归于祂的神性；并且一方面说\Quote{你们把拉匝禄放在哪里？}\BibleRef{若 11:34}属于我们，而大声呼喊\Quote{拉匝禄，出来！}以及使已死四天的人复活，则超乎我们的本性；另一方面，虽然\Quote{祂在痛苦中、祂被钉十字架、祂被埋葬}属于帏幔，但\Quote{祂有信心、祂复活、祂升天}属于内藏的宝库；然后他们指责我们引入两种分开或冲突的本性，并分割那超自然且奇妙的结合。他们要么不该做他们指责我们的事，要么不该指责我们做他们自己做的事；至少如果他们决心保持一致，而不同时提出他们自己和对手的原则的话。这就是他们缺乏理性的程度；它与自身和真理都如此冲突，以致于当他们与自己不一致时，既没有意识到也不感到羞愧。现在，如果有人以为我们是自愿而非出于勉强写这一切，以为我们在劝阻合一，而非尽最大努力促进它，那么让他知道他大错特错，对我们的渴望猜测得完全不对，因为在我们眼中，没有什么比和平更珍贵，过去也从未有过，事实本身证明了这一点；尽管他们的行动和针对我们的吵闹完全排除了同心合意。\par

% structure: p0030 source=h2 target=subsection reason=relative-heading-rank; base=h2

\subsection{书信一七五。致\Person{Olympius}。}

即使白发苍苍也还有可学之处；老年，似乎也不能在所有方面都信赖其智慧。至少我自己，比任何人更了解阿波利拿里派的思想和异端，看出他们的愚昧是不可容忍的；然而，我以为自己能以忍耐驯服他们，逐渐软化他们，便让自己的希望促使我渴望达到这一目标。但是，看来我忽略了这样一个事实：我反而使他们变得更坏，并以不合时宜的哲学损害了教会。因为温和并不能使恶人感到羞愧。如今，如果我能亲自将这些教给你，我当然不会犹豫，甚至会不辞旅途劳顿，俯伏在你的尊驾面前。但由于我的病情已使我过于虚弱，并且遵照医生们的建议，我必须尝试\Place{克桑克萨里斯}的温泉浴，因此，我寄去一封信代表我。这些邪恶和彻底堕落的人，竟敢在他们所有其他恶行之外，要么召集某些主教，要么滥用某些主教的段落（我不准备确切说明是哪一种），这些主教是由东西方全教会会议剥夺职务的；并且，他们违反所有帝国法令和你的命令，将主教的名称授予他们自己那伙不信和欺骗的团伙中的某个人；我相信，他们之所以如此胆大妄为，没有比我的极大虚弱更鼓励他们的了；因为我必须提到这一点。如果这是可以容忍的，你的尊驾将会容忍它，而我也将像以前多次那样忍受。但如果它是严重的，并且不能被我们最尊贵的皇帝们所容忍，那么，请惩处所发生的事——尽管要比这种疯狂应得的惩罚更温和些。\par

\SourceRecord{Translated by Charles Gordon Browne and James Edward Swallow.；Nicene and Post-Nicene Fathers, Second Series；Vol. 7.；Edited by Philip Schaff and Henry Wace.；Buffalo, NY: Christian Literature Publishing Co.,；1894.；<http://www.newadvent.org/fathers/3103a.htm>.}

% structure: p0001 source=h1 target=section reason=page-title

\section{\WorkTitle{Letters (Division II)}}

\Emph{与凯撒勒雅总主教圣巴西略之通信}\par

% structure: p0003 source=h2 target=subsection reason=relative-heading-rank; base=h2

\subsection{第一书　致同修巴西略}

（约公元357或358年；答复现已不存的一封信。）\par

我承认，我未能遵守诺言。早在雅典时，在我们彼此友谊与亲密交往之际（我找不到更好的词来形容），我曾承诺要与你共度哲学生活。但我未能履行诺言，并非出于我自己的意愿，而是因为一条法律战胜了另一条法律：我说的是那要求我们孝敬父母的法律，胜过了我们友谊与交往的法律。然而，你若接受这个提议，我不会完全辜负你。我将与你共度一半时间，另一半你与我共度，这样我们就共同拥有全部，我们的友谊就可平等；这样安排，既不使我的父母悲伤，我又能得着你。\par

% structure: p0006 source=h2 target=subsection reason=relative-heading-rank; base=h2

\subsection{第二书}

（约同时期写，答复另一封现已不存的书信。）\par

我可不喜欢你拿台伯里纳及其泥泞和冬天开玩笑，我的朋友，你是那样一尘不染，踮着脚尖走路，踏着平原。你长了翅膀，高高飞翔，像阿巴里斯之箭一样飞行，为的是，尽管你是卡帕多细亚人，却要逃离卡帕多细亚。我们亏待了你吗？当你面色苍白、气喘吁吁、测量着太阳时，我们却皮肤光滑、吃得饱饱的、不觉得拥挤？可这就是你的状况。你奢侈富有，还去赶集。我不赞成这样。那么，要么你停止指责我们的泥泞（因为既不是你建了你的城，也不是我们创造了我们的冬天），要么为了我们的泥泞，我们要用你的小贩以及城市里滋生的其他麻烦来对抗你。\par

% structure: p0009 source=h2 target=subsection reason=relative-heading-rank; base=h2

\subsection{第四书}

（答复巴西略第十四书，约361年。）\par

你可以嘲笑并拆解我的事情，无论是开玩笑还是当真。这无关紧要；只管笑吧，尽情享受文化，享受我的友谊。你所说的一切，对我来说都是愉快的，无论它是什么，看起来如何。因为我认为你嘲弄这里的事情，并非为了嘲弄本身，而是为了吸引我到你自己那里，如果我对你有所了解的话；就像人们阻塞水流，以便将其引入另一水道一样。你的言语总是给我这样的感觉。\par

至于我，我将赞赏你的本都和你本都的黑暗，你那值得流放的居所，你头顶的山丘，考验你信德的野兽，以及你那隐蔽之所……或者我应该说，你那老鼠洞，却冠以\InnerQuote{思想之所}、\InnerQuote{隐修院}、\InnerQuote{学院}等堂皇名号；你的灌木丛和陡峭山峰的花环，愿你被它们环绕，而非加冕；你有限的空气；你渴望却只能像透过烟囱看到的太阳，哦，本都的无日者，不仅像有些人所说的被判处六个月的黑夜，而且连你生命的一部分都没有脱离阴影，你的整个生命都是一个漫长的黑夜，真正的死亡阴影，用圣经的话来说。还有赞赏你那狭窄艰苦的道路，引向……我不知道是天国还是冥府，但为了你的缘故，我希望是天国。至于中间的地区，你希望怎样？我是否要谎称它为伊甸园，分为四道河流，灌溉世界？还是干涸无水的旷野（只是哪一位梅瑟会来用杖从岩石中引出水面来驯服它呢）？因为所有避开岩石的地方都是沟壑；不是沟壑的地方则是荆棘丛；荆棘之上的则是悬崖；其上的路是陡峭的，两边倾斜，锻炼着行人的心智，需要体操练习来确保安全。河流咆哮着冲下，对你来说，它像安菲波利斯的斯特律蒙河一样安静，但里面的鱼不比石头多，它不流入湖泊，而是冲入深渊，哦，我那位夸大其词的朋友和新名字的发明者。因为它又大又可怕，淹没了上面居民的圣咏；如同尼罗河的瀑布和急流，它日夜咆哮着压倒你。它湍急无渡；只有这一点小仁慈：当洪水和冬天的暴风雨使它狂暴时，它不会冲走你的住所。这就是我对那些福岛和你们幸福之人的看法。你不该赞赏那些新月形的弯曲，它们扼杀而非切断你们高地的可通行部分，以及悬在你们头顶的山脊带，使你们的生活像坦塔罗斯一样；还有通风的微风和地表的通气孔，当你们勇气不足时，它们使你们振作；还有你们那唱歌（但只唱饥饿）和飞翔（但只在沙漠中）的音乐之鸟。你说没有人来访，除非为了打猎；你还可以补充说，除非为了来看你们的尸体。这封信也许太长了，但比起一部喜剧又太短。如果你能善意地接受我的玩笑，那很好；不然，我会再寄些给你。\par

% structure: p0013 source=h2 target=subsection reason=relative-heading-rank; base=h2

\subsection{第五书}

（约公元361年。）\par

既然你善意地接受了我的玩笑，我把余下的寄给你。我的序曲来自荷马。\par

现在且改变你的主题，\newline{}歌唱内在的装饰。\par

—— \WorkTitle{奥德赛}卷八第492行。\par

你那没有屋顶、没有门的小屋，你那无火无烟的炉灶，你那被火烤干的墙壁——免得我们被泥点打到——就像坦塔罗斯在汪洋中焦渴，还有那场可悲的宴席，没有食物可吃；我们是从卡帕多细亚被邀请来的，不是去赴食忘忧果者式的清贫，而是去赴阿尔西诺斯的筵席——我们这些年轻可怜的沉船幸存者。因为我记得那些面包和那汤（姑且这么叫吧），是的，我还会记得它们，还有我那可怜的牙齿，在你的硬面包上打滑，然后又绷紧，仿佛从泥潭里拔出来。你自己会把这一切提升到更惨烈的悲剧高度，因为你已经学会了通过自己的苦难夸夸其谈……要不是那位穷人的伟大支持者——我说的是你的母亲——迅速解救了我们，她像暴风雨中颠簸之人的港湾一样及时出现，我们早就死了，人们会为我们在本都的信心感到可怜多于钦佩。我怎能不提那个不是花园、没有蔬菜的花园，还有那个我们从屋里清除出来、又把它填满（花园）的奥革阿斯的粪堆，当我们拉着那座山一般的货车时，我是葡萄采收者，你是勇士，用我们的脖子和手，至今还留着我们劳作的痕迹。\Quote{大地啊，太阳啊，空气啊，美德啊}（因为我将稍作悲剧性的语调），不是为了架设赫勒斯滂的桥梁，而是为了削平一座悬崖。如果你对这些事情的提及不感到恼火，我也不恼火；但如果你恼火，我对现实就更恼火。其余的我就略过了，出于对其他人的尊重，我从他们那里得到了许多快乐。\par

% structure: p0019 source=h2 target=subsection reason=relative-heading-rank; base=h2

\subsection{第六封信。}

（约写于同时，语气更为严肃。）\par

我之前关于我们在本都的那段时光是戏言，并非认真的；而现在我写的是非常认真的。但愿有人能将我置于往昔那些日子之月，\BibleRef{约 29:2}——那时我与你在艰苦生活中一同欢乐；因为自愿的痛苦比非自愿的快乐更宝贵。但愿有人能归还我那些咏唱、守夜，以及与天主在祈祷中的同住，还有那种可以说是无形无体的生活。但愿有那种兄弟间的亲密与同心，他们被你神圣化了，得到升华；但愿有那种通过成文的规则与规范而确保的修德竞赛与激励；但愿有在圣经中的热诚劳作，以及在圣神的引导下我们从其中发现的光明。或者，若我可以说些更小更轻的事，但愿有日课与经验；但愿有拾柴与切石；但愿有那棵金色的悬铃木，比薛西斯的更珍贵，坐在它下面的不是一位被奢华削弱的国王，而是一位被艰苦生活磨砺的隐修士，我栽种了，阿颇罗（即尊敬的您）浇灌了，\BibleRef{格前 3:6}但天主使它生长，成为我们的荣耀，好让我的勤勉在你们中间留下纪念，正如我们在约柜中读到并相信的，亚郎发过芽的棍杖。\BibleRef{户 17:8}渴望这一切很容易，但达到却不容易。但你到我这里来，与我一同修德，与我合作，以你的祈祷帮助我保持我们曾共同获得的益处，免得我像日暮时的影子一样渐渐消逝。我宁愿呼吸你胜过空气，只有与你在一起时才活着——或是实际与你同在，或是因你的相似而虚拟地在你不在时与你同在。\par

% structure: p0022 source=h2 target=subsection reason=relative-heading-rank; base=h2

\subsection{第八封信。}

（写于圣巴西略被祝圣为司铎后不久，大概在公元362年底。）\par

我赞成你信的开头；但你的什么我不赞成呢？你被证实写得完全像我；因为我也是被迫进入司铎之列的，因为我从未渴望过它。我们彼此是可靠见证人，证明我们热爱谦卑低调的哲学。但也许这没有发生更好，或我不知道该说什么，因为我不明白圣神的旨意。但既然事情已经发生，我们必须承受，至少在我看来是这样；尤其是当我们考虑到时代，这些时代给我们带来了如此多的异端舌辩，我们绝不能辜负那些如此信赖我们的人对我们的期望，也不能辜负我们自己的生命。\par

% structure: p0025 source=h2 target=subsection reason=relative-heading-rank; base=h2

\subsection{第十九封信。}

（这封信应与所提及的写给凯撒勒雅主教欧瑟伯的三封信一起阅读。背景情况见总绪论，§1，第194页。）\par

现在是需要谨慎和忍耐的时候，我们不应让任何人显得比自己更有勇气，也不应让我们所有的劳作和辛苦瞬间化为乌有。我为何写这些？因为什么？我们的主教欧瑟伯——非常蒙天主喜悦（因为今后我们必须这样想、这样写他）——非常倾向于与我们的和解与友谊；如同火使铁变软，时间使他变软了；我想，一封恳求和邀请的信将从他那里寄给你，正如他向我暗示的那样，许多熟悉他事务的人也向我保证。那么，让我们先走一步，要么去他那里，要么写信给他；或者先写信，然后再去；以免到时候我们被打败而感到羞愧，而原本有能力获得胜利——通过被体面地、明智地击败，这正是许多人向我们要求的。请听我的劝，来吧；既是因为这个原因，也是因为时势艰难；因为异端者的阴谋正在攻击教会；他们中有些人现在就在这里，正困扰我们；其他人，传闻说，正在到来；有理由担心真理之言会被冲走，除非很快兴起一个贝匝肋耳的精神，他是这类论点和教义的总建筑师。如果你认为我也应该去，和你一起住、一起旅行，我甚至不会拒绝这样做。\par

（我们在此插入上述三封致欧瑟伯的信件，因它们与上文紧密相连，不至于显得突兀。）\par

% structure: p0029 source=h2 target=subsection reason=relative-heading-rank; base=h2

\subsection{第十六封信。致凯撒勒雅主教欧瑟伯。}

既然我现在致书的对象是一位不爱谎言的人，而且据我所知，在他人面前识破谎言时，无论它如何巧妙地缠绕成各种迷宫，他都是最敏锐的；此外，就我自身而言——尽管我违心地说——我也不喜欢诡计，这既出于我的天性，也因天主的圣言如此造就了我。因此，我写下心中所现，并恳请您宽恕我的直言不讳；否则，您将因剥夺我的自由、强迫我将悲伤的痛苦像某种隐秘的恶疾般压抑在内心而损害真理。我欣喜自己能蒙受您的敬重（因为我正如某人先前所说，是一个人），并被召往主教会议和灵修会议。但我为我的最可敬的兄弟巴西略所遭受且仍在受您阁下轻视之事而忧虑；因为我选定他作为我生活、言谈及最高哲思的同伴，他至今仍是如此，我从未后悔对他的判断。这样谈论他已算是温和，以免因赞美他而显得自我夸耀。然而，我认为您通过尊崇我而羞辱他，行为恰似一人用一只手抚摸他人头颅，另一只手却击打其面颊；或一边掘毁房屋根基，一边粉刷墙壁并装饰外表。因此，若您肯听从我，这便是您当行之事，而我恳请被倾听，因为这是公义。若您对他给予应有重视，他也会同样待您。我将如影子随形一般跟随他，因我微不足道且倾向于和平。因为我不至于卑劣到在其余方面愿意追求哲学并归于上品，却忽略我们全部教导的终向——即爱德；尤其关乎一位司铎，且是品性如此崇高者，并且据我所知，在所有相识中，他在生活、教义和行为上皆是最优秀的。因为我的痛苦不会掩盖真理。\par

% structure: p0031 source=h2 target=subsection reason=relative-heading-rank; base=h2

\subsection{第十七封信。致凯撒勒雅总主教欧瑟伯。}

我写信并非出于傲慢之心，正如您埋怨我的信那样，而是出于一种灵修与哲学的心态，并且是合宜的——除非这也冒犯了\Quote{您的雄辩的额我略}。因为尽管您在地位上高于我，但您会赐予我些许自由和坦率直言的权利。因此，请您对我宽厚些。但若您将我的信视为出自仆人之手，且来自一个连正视您面孔都不配的人，那么这次我愿接受您的鞭笞，甚至不流一滴泪。您也会因此责怪我吗？这更适合于任何其他人，而非您阁下。因为高尚灵魂之人更愿接受朋友的自由，而非敌人的谄媚。\par

% structure: p0033 source=h2 target=subsection reason=relative-heading-rank; base=h2

\subsection{第十八封信。致凯撒勒雅的欧瑟伯。}

我从未对您阁下怀有卑劣之心；请勿判我有罪。但在容许自己稍许自由与勇气、仅仅为了舒缓和医治我的悲伤之后，我随即俯首服从，甘愿顺服于教规。我还能做什么呢？我既认识您，也认识圣神之律。然而，即使我的情感曾如此卑劣和鄙下，现今之时也不容许此类情绪，何况那些冲向教会的野兽，以及您自身为教会而战时那般纯洁而真切的勇气与刚毅。因此，若您愿意，我将来临，与您一同祈祷并争战，服侍您，如同欢呼的孩童，以我的劝勉鼓舞这位高贵的斗士。\par

% structure: p0035 source=h2 target=subsection reason=relative-heading-rank; base=h2

\subsection{第四十封信。致大巴西略。}

\EditorialNote{约在公元370年中。欧瑟伯去世后，巴西略似乎有意让其朋友额我略继承空缺的总主教之位；于是他写信给额我略，未提总主教之死，而表示自己病危，请他前来凯撒勒雅。额我略深为悲痛，立即动身，但在途中得知巴西略身体如常，而教省的主教们正聚集凯撒勒雅选举总主教。他立刻识破了这一计谋；并认为巴西略希望自己在都城出现，以借助其影响促成他（巴西略）的当选，因而写了下面这封愤怒的信。不过，额我略及其父亲都感到无人比巴西略更适合接掌空位；于是额我略以父亲的名义写了后面三封信，分别致凯撒勒雅民众、出席主教会议的主教们以及萨摩萨塔主教欧瑟伯。}\par

若我说出一些奇特的、前人未曾说过的话，请勿惊讶。我认为你有稳健、可靠且意志坚定的名声，但在你的许多计划和行动中，却显得过于单纯而非稳健。因为那远离邪恶的，也相应地不易猜疑邪恶，正如刚才所发生的事所示。你在我被召集到大都会之际，正值为选举主教而召开会议，而你的借口十分得体且貌似合理。你假装病重，已奄奄一息，渴望见我并向我作最后的告别；我不知道其目的何在，甚至不知道我的到场能对这事有何影响。我为所发生的事深感悲痛地出发；因为对我来说，有什么比你的生命更宝贵，或比你的离世更令人悲痛呢？我泪如泉涌，放声哀哭，第一次感到自己失去了哲学家的镇定。我哪一样送葬的事没有做到？但当我发现主教们正在城中集会时，我立刻停止了行程；首先我奇怪的是，你竟没有觉察到什么是合宜的，或防范人们的口舌——它们如此惯于诽谤坦诚的人；其次，你竟不认为同样的事对我和对你是合适的，尽管我们的生活、规矩和一切对我们两人都是共通的，我们自始就由天主如此紧密地结合在一起。第三（因为我也必须说这一点），我奇怪你是否记得，这样的任命更适于更虔诚的人，而非更有权势的人，也非最得民心的人。因这些缘故，我便后退并迟疑了。如今，若你与我想法相同，就作出这个决定吧：避免这些公共的纷争和邪恶的猜疑。等事情安定、时间允许时，我将拜见你的尊驾，并将向你提出更多更严厉的责备。\par

% structure: p0038 source=h2 target=subsection reason=relative-heading-rank; base=h2

\subsection{第四十一书。以他父亲名义，致\Place{凯撒利亚}民众。}

我是一个小小的牧人，管理着一小群羊，且在圣神的仆人中是最小的一个。但恩宠并不狭隘，也不受地点限制。因此，即使渺小的人——尤其是当主题如此重要，且人人关心时——也应有言论自由，甚至可与那些白发苍苍、或许比常人更有智慧的人一同商议。你们所商议的并非寻常小事，而是关乎根据你们的决定，公共利益必然受到促进或损害的事。因为我们的话题是教会——基督为之而死——以及将要引领教会走向天主的牧者。因为身体的光是眼睛\BibleRef{玛 6:22}，正如我们所听见的；不仅是看得见并被看见的肉眼，也是那以属灵方式观看和被人观看的心眼。但教会的光是主教，即使我们不写，对你们也是显而易见的。因此，正如身体的直行或弯曲取决于眼睛的清晰或昏暗，教会也必然因其首领的行为而分担危险或安全。所以你们必须为整个教会——基督的身体——着想，但更应为你们自己的教会着想，它从起初就是几乎众教会的母亲，现在仍是；所有的共同体都注视着它，如同围绕圆心所画的圆，这不仅是因为它自古以来向众人宣告的正统信仰，也是因为天主显然赐予它的合一的恩宠。你们也邀请我们参与你们对这件事的讨论，这做得正确且符合教规。然而我们为年龄和病弱所困，如果我们能藉着圣神所赐的力量出席（对于相信的人，没有不可能的事），那么对公共福祉最好，对我们自己最愉快，以便我们能有所贡献于你们，并自己也有份于祝福；但如果我因软弱而不能前往，我无论如何将提供缺席者所能提供的一切。\par

我相信你们中另有其他人堪当首席之职，既因你们城市的伟大，也因它过去曾由如此卓越的伟人治理得如此出色；但在你们中间有一个人，我无法把任何人置于他之上，即我们蒙天主眷爱的儿子，司铎\Person{巴西略}（我在天主面前作见证）；他生活纯洁，言语正直，独自——或几乎是独自——在两方面都有资格抵挡现今的时代和异端者流行的空谈。我写信给司铎和隐修会士们，也给显贵和议员们，以及全体民众。如果你们赞同我的意见，而我的投票既如此公正合理，又是在天主的助佑下作出的，那么我心神与你们同在，并将继续同在；或者不如说，我已着手此事，并因圣神而勇敢。但如果你们不同意我，而另作决定，如果此事将由派系和关系来决定，如果群众的插手再次扰乱你们投票的诚意，那么你们就做你们喜欢的事吧——我将留在家中。\par

% structure: p0041 source=h2 target=subsection reason=relative-heading-rank; base=h2

\subsection{第四十三书。致主教们。}

\EditorialNote{（同省的主教们已通知年长的\Person{额我略}他们的主教会议，但没有提及日期或目的，也未邀请他参加——大概因为他们知道他会大力支持\Person{巴西略}的选举，而他们对巴西略不利。因此圣\Person{额我略}以他父亲的名义写了以下信函。）}\par

你是何等甘饴、何等仁慈，又是何等充满爱德。你邀请我前往大都会，因为我想，你将要就一位主教的事进行商议。我从你信中得知了这些，尽管你既未告诉我是否要我出席，也未说明缘由或时间，只是突然告知我你已动身，仿佛决意不尊重我，且不愿我参与你的商议，反而阻碍我前来，免得你甚至违背我的意愿与我相见。这是你行事的方式，我愿忍受此侮辱，但我将向你陈明我的看法与感受。各色人等会提出不同的候选人，各依其爱好与私利，正如在这种时刻常见的那样。但我不能偏爱任何人，因我的良心不允许我这样做，除非是我亲爱的儿子、同司铎巴西略。因为在我所有相识中，我找得出谁比他生活更经得起考验，言词更有力量，或全然更具美德的光彩？但若你以他身体虚弱为由反对，我回答说，我们选的是教师，而非运动员。同时，在此事上，我们看到那加强并扶持软弱者（倘若他们如此）的天主的能力。若你接纳此票，我将前来参与，或心神，或身体。但若你走向一个既定的结论，且派系将凌驾正义，我将因被忽略而欢喜。事工必属于你，但请为我祈祷。\par

% structure: p0044 source=h2 target=subsection reason=relative-heading-rank; base=h2

\subsection{书信第四十二。致\Place{萨莫萨塔}主教\Person{欧瑟伯}。}

\EditorialNote{当时似乎仍有阴谋和党派精神得胜的可能，因此两位\Person{额我略}决定邀请\Place{萨莫萨塔}的\Person{欧瑟伯}的援助，尽管他不属于该省。他一直是正统信仰对抗亚略派皇帝\Person{瓦伦斯}的显赫捍卫者，两位\Person{额我略}对他的出席主教会议寄予厚望。他回应了他们的请求，跋涉三百英里极为艰难的路程，以将他的影响力投入到他们所关切的志业中。然而，他看出，年迈的\Place{纳齐盎}主教尽管年事已高、身体衰弱，也必须做出努力，于是他说服了他前往。结果尽如人意：\Person{巴西略}获得一致选举通过。\Person{圣额我略}后来以自己的名义写信感谢他，此信将在后面出现。}\par

但愿我有鸽子般的翅膀，或我的老年能恢复青春，使我能来到你的爱德前，满足我见你的渴望，并向你倾吐我灵魂的忧苦，在你的安慰中找到对我痛苦的慰藉。因为自从真福主教\Person{欧瑟伯}去世以来，我颇为担心，唯恐那些从前为我们的\Place{大都会}设下陷阱、欲以异端的莠子充满它的人，现在会抓住机会，用他们邪恶的教导拔除那在人心田中付出如此辛劳播种的虔敬，并撕裂其合一，正如他们在许多教会中所行的那样。我一收到圣职人员的来信，请求我在此情势下不要忘记他们，我便环顾四周，记起了你的爱德、你的正确信仰以及你对天主教会常怀的热忱；因此，我派遣我所爱的执事兼助手\Person{欧斯塔丢}前去提醒你的尊威，并恳求你，除了你为教会已付出的一切辛劳之外，前来与我会面，既以你的到来使我的老年得到慰藉，又在那如此著名的正统教会中确立虔敬，与她一同（若我们被堪当与你共享善工）按主的旨意赐予她一位能治理祂子民的牧者。因为在我们眼前有一个人，你并非不认识他；若我们被允许得到他，我知道我们将获得向天主的大胆信心，并将给那呼求我们援助的民众带来极大的益处。我再次恳求你，放下一切迟延，在冬季恶劣天气来临之前前来我们这里。\par

% structure: p0047 source=h2 target=subsection reason=relative-heading-rank; base=h2

\subsection{书信第四十五。致\Person{巴西略}。}

\EditorialNote{祝圣之后，人人都以为\Person{额我略}会立即加入他的朋友；\Person{巴西略}自己也极希望得到他的协助。但\Person{额我略}认为，最好克制自己见朋友的愿望，等嫉妒情绪有时间平息下来。因此，他写了以下信件，解释他在此时不前往的理由。}\par

当我得知你被安置在高高的宝座上，且圣神已得胜，将那灯台上的蜡烛——它即使在先前也非暗淡之光——显明出来时，我承认，我欢喜。我为何不喜呢？因为我看见教会的公团处于可怜的境地，需要如此一位引导之手。然而，我并未立刻跑向你，也不会跑向你，即使你自己邀请我。首先，为谨慎维护你的尊位，免得你显得是在不良品味和急躁情绪下招聚党羽，正如你的诽谤者会说的那样；其次，为使自己获得稳健和超越恶意的名声。那么，你何时会来呢？也许你会问，你还要延迟多久？直到天主吩咐我，直到当前敌意和诽谤的阴影过去。因为我深知，那些癞病人不会长久阻挡我们的达味进入耶路撒冷。\par

% structure: p0050 source=h2 target=subsection reason=relative-heading-rank; base=h2

\subsection{书信第四十六。致\Person{巴西略}。}

\EditorialNote{新任总主教似乎对\Person{额我略}上一封信中给出的理由并不满意；因此后者再次写信。}\par

你的事务怎能对我只是拾取葡萄的零碎呢，亲爱的神圣朋友？\par

\Quote{有什么话从你牙齿的栅栏中逃了出来？}或者，若我也可稍微大胆一点，你怎么敢说出这样的话？你的心如何使它发出，你的墨如何写下，你的纸如何承受，哦，讲演、\Place{雅典}、德行和文学的劳作！你所写的几乎使我写出悲剧来。难道你不认识我，或不认识你自己，你这世界的眼、伟大的声音、号角和学问的宫殿？你的事在\Person{额我略}看来是琐事？那么，在地上还有什么能让人钦佩，如果\Person{额我略}不钦佩你？在季节中有一个春天，在星辰中有一个太阳，有一个天空包容万物；同样，在所有事物中，你的声音是独一无二的——如果我能判断这类事，并且没有被我的情感所欺骗——而我认为事实并非如此。但若因为你责备我，说我没有按你的价值来评价你，那么你也必须责备全人类；因为没有人曾经或将要足够钦佩你，除非是你自己和你自己的口才——至少若是可能赞颂自己，且若这是我们的言谈习惯。但若你控告我轻视你，为何不更是控告我发疯？或者你恼怒是因为我行事像一个哲人？请允许我说，这一点，也只有这一点，甚至高于你的交谈。\par

% structure: p0054 source=h2 target=subsection reason=relative-heading-rank; base=h2

\subsection{第四十七封书信。致\Person{巴西略}。}

（公元372年，\Place{卡帕多细雅}民政省分为两省，随之而来的是教会内的纷争。新\Place{卡帕多细雅}第二省民政首府\Place{提雅纳}的主教\Person{安提摩斯}主张，教会分区必须与民政分区一致，因此声称国家纯粹民政的行动\OriginalTerm{当然}{ipso facto}已将他提升为新省都主教的尊位；这一主张得到了该地区主教们的支持，他们通常对这位伟大的总主教并无好感。以下三封信便与这一争端有关。）\par

我听说你正因这一新事物而受困扰，并因当权者中某些诡辩且并非不寻常的多管闲事而烦恼；这并不奇怪。因为我并非不知他们的嫉妒，也不知你周围的许多人正在利用你来谋取自身的利益，并点燃卑劣的火花。我不怕看见你因困扰而失去哲人的心境，或做出任何有损于你和我之事。相反，我认为现在正是我的\Person{巴西略}将为人所知的时刻，你一生所积存的哲思将显现出来，它将如高浪般克服凌辱；你将屹立不动，而他人则受扰动。你若认为好，我就亲自前来，或许能凭我的劝言给你些许帮助（若大海需要水，你确实需要劝言！）；但无论如何，我将获益，并因分担一部分凌辱而学到哲思。\par

% structure: p0057 source=h2 target=subsection reason=relative-heading-rank; base=h2

\subsection{第四十八封书信。致\Person{巴西略}。}

（在上述事件之后不久，\Person{巴西略}决定在争议省份中设立若干新的主教区以加强己方力量，其中之一为\Place{萨西玛}，他将\Person{额我略}祝圣为那里主教，这大大违背了后者的意愿；\Person{额我略}感到受了亏待，并不掩饰其不情愿。见\WorkTitle{总导论}第195页。）\par

请不要再说我缺乏教养、粗鲁无礼、不友善，甚至不配活着，因为我敢于意识到自己所受到的待遇。你自己也会承认，我在其他方面并没有做错；我的良心也不指责我曾在大小事上对你刻薄，并且我希望永远不会如此。我只知道我看到自己受到了欺骗——虽然太迟了，但我确实看到了——而我将责任归咎于你的宝座，因为它突然使你高升于自己之上。我厌倦了因你的过错而受到责备，也厌倦了向那些了解我们过去和现在关系的人为这些过错辩解。因为在我必须忍受的一切中，最可笑或最可怜的是，同一个人既要承受委屈，又要承担责备，而这正是我目前的处境。不同的人根据各自的喜好、性情或对我情绪的深浅，用不同的理由责备我；但最善意的人也指责我傲慢和轻蔑，他们在利用了我之后，就把我抛在一边，如同最无用的器皿，或那些建拱门时用的框架，建筑完成后就被拆下丢弃。就让他们去吧，随他们说什么；没有人会限制他们的言论自由。而你，作为我的报酬，请偿还那些你为对付诽谤者而编造的空洞祝福和虚假希望——那些诽谤者指责你以尊重的名义侮辱我，好像我轻浮易欺，轻易被这种对待所蒙骗。现在我要坦诚说出我的心境，你不可对我生气。因为我将告诉你，在受苦的那一刻我说过的话，并非出于愤怒，也不是因惊愕而失去理智或不知所云。我不会拿起武器，也不会学习我以前在时机看似更合适时（当时人人武装，处于狂热之中；你知道弱者之病）未曾学会的战术；我也不会赤手空拳、不习战事地去面对好战的\Person{安提摩}，尽管他是个不合时宜的武士，而我这样反而更容易受伤。如果你愿意，你自己与他战斗吧（因为需要往往使弱者也成为战士），或者另找他人，在他夺取你的骡子、把守隘口、像昔日的\Person{阿玛肋克}阻挡以色列的道路时去战斗。请首先赐我宁静。我为何要为小猪和鸡鸭（况且不是自己的）战斗，仿佛为了灵魂和教规？我为何要剥夺大都会那著名的\Place{萨西玛}，或揭露你心底的秘密，而我本应协助掩盖它？那么，你要显示男子气概，刚强起来，把各方都引向你自己的结论，如同河流将冬天的洪水卷入，不顾及友情或亲密的善，也不顾及如此行为会带给你的名声。只把自己交托给圣神。我从你的友谊中只能得到这一点：学会不再信任朋友，也不把任何事物看得比天主更宝贵。\par

% structure: p0060 source=h2 target=subsection reason=relative-heading-rank; base=h2

\subsection{书信49。致\Person{巴西略}。（宁静的赞美。）}

你指责我懒惰和闲散，因为我没有接受你的\Place{萨西玛}，没有像一位主教那样忙碌，没有像投到狗群中的骨头一样挑起你们彼此争斗。我最大的事务始终是远离事务。为了让你了解我的一个优点，我如此珍视远离事务，以至于我认为自己甚至可以成为所有人在这类宽宏大量上的标准；只要人人都效仿我，教会就不会有纷扰，那被每个人当作私斗武器的信仰也不会被撕碎。\par

% structure: p0062 source=h2 target=subsection reason=relative-heading-rank; base=h2

\subsection{书信50。致\Person{巴西略}。}

（似乎是应\Person{安提摩}的请求，\Person{圣额我略}曾致函\Person{圣巴西略}（该信现已不存），提议两位对立的\OriginalTerm{总主教}{Metropolitans}举行一次会议。\Person{圣巴西略}对这个善意的提议感到不快，并写了一封生硬的信给\Person{圣额我略}，以下是对该信的回复。）\par

你在书信中跳脱如马驹，何等热烈。我并不奇怪，当你刚成为光荣的产业，你就想向我显示你所发现的光荣是什么，好使自己更显威严，就像那些描绘季节的画家一样。但是，要解释关于主教们的全部事情以及那封使你烦恼的信，我的出发点、我走了多远、在哪里停下，在我看来，对于一封信来说太长了，而且这个主题与其说是辩护，不如说是历史。我简洁地向你解释：最尊贵的\Person{安提摩}与几位主教来到我们这里，不知是来拜访我的父亲（至少这是借口），还是来做他所做的事。他在许多方面、许多话题上试探我：教区、\Place{萨西玛}的沼泽地、我的祝圣……恭维、质问、威胁、恳求、责备、赞扬、在自己周围画圈，仿佛我只应看他和他那新的都主教区，因为它是更大的。我说，你为什么画线把我们的城市也包括进去？因为我们也认为我们的教会确实是教会之母，而且从古时就是如此。最后，他没有达到目的，气呼呼地离开了，并以\OriginalTerm{巴西略主义}{Basilism}责备我，仿佛那是一种\OriginalTerm{腓力主义}{Philipism}。你认为我在这件事上做了错事吗？现在看看我的这封信——你说是侮辱了你的那封。他们拟定了一份主教会议的召集令给我；当我拒绝并说这是侮辱时，他们转而要求通过我来邀请你商讨这些事情。我答应了，以防止他们的第一个计划得以实施；将整个事情交在你手中，如果你选择召集他们，以及在哪里、何时召集。如果我在这件事上没有伤害你，请告诉我哪里有可能伤害。如果你要从我这儿得知，我会把\Person{安提摩}在入侵沼泽地之后寄给我的信读给你听，尽管我禁止和威胁，他侮辱和谩骂我，仿佛在为我失败而唱凯歌。为什么我要为了你而冒犯他，同时又使你不快，好像我在讨好他？亲爱的朋友，你首先应该知道这一点；即使如此，你也不该侮辱我——只因为我是一名司铎。但如果你非常倾向于炫耀和争吵，并以我上级的口吻说话——如同都主教对一个微不足道的辅理主教，甚至对一个没有教区的主教——那么我也有点骄傲来对抗你。这对任何人来说都很容易，也许是最合适的做法。\par

% structure: p0065 source=h2 target=subsection reason=relative-heading-rank; base=h2

\subsection{书信五十八。致巴西略。}

\EditorialNote{有人在额我略面前攻击巴西略关于天主圣神的天主性的正统信仰；在这封信中，他向友人叙述了自己辩护的方式。不幸的是，巴西略对这封信并不满意，认为它意在假借友好的同情感传达责备。}\par

从一开始，我就把你——现在仍把你——当作我生活的向导和信德的导师，以及一切可称颂之事的主宰；若有人赞美你的功德，他完全与我同在，甚至在我之后，因你的虔诚如此超越我，我完全属于你。这并不奇怪，因为交往愈久，经验愈深；经验愈丰富，见证愈完美。若我在生活中有所获益，那都是来自与你的友谊和交往。这就是我对这些事的态度，并希望永远如此。我现在所写的，是出于不情愿，但我仍然写了。不要生我的气，否则我也会很生气，如果你不相信我是出于善意说和写的。\par

许多人谴责我们在信德上不坚定；我是指那些认为我们完全一致的人——他们想得对。有些人公开指控我们异端，另一些人指控我们懦弱；异端——那些认为我们的言论不纯正的人；懦弱——那些责备我们缄默的人。我不必转述别人说什么；我将告诉你最近发生的事。\par

这里有一个聚会，许多杰出的朋友都在场，其中包括一位穿着表明虔诚的名号和服装的人（即一位隐修士）。他们尚未开始饮酒，但在谈论我们，正如这类聚会上常有的那样，并且把我们而不是其他事情作为谈话的主题。他们赞赏你的一切，并把我引为同样哲学的宣扬者；他们谈到我们的友谊、雅典，以及我们在所有观点和感情上的一致。我们那位哲学家对此感到恼火。\Quote{诸位，这是什么意思？}他大声喊道，\Quote{你们是多么说谎和谄媚的人。如果你们愿意，可以因其他理由赞扬这些人，我不会反对；但最重要的一点——他们的正统信仰——我不能让步。巴西略和额我略受到虚假的赞扬；前者是因为他的言语背叛了信仰，后者是因为他的容忍助长了背叛。}\par

我说，这是什么意思，你这虚妄的人，新\Person{达堂}和\Person{阿彼兰}，愚妄之极？你从哪里来，给我们制定法律？你凭什么自以为是如此重大之事的审判官？\Quote{我刚从\Person{欧普萨基}殉道者的庆节回来，}他回答说（事实确实如此），\Quote{在那里我听见伟大的巴西略最美妙、最完美地讲述了圣父和圣子的天主性，几乎无人能及；但他对圣神却一带而过。}他还加上了一种关于河流的比喻，说河流经过岩石，冲蚀出沙土。\Quote{至于你，我的好先生，}他看着我说，\Quote{你现在确实公开表达了对圣神的天主性的看法，}他提到我在一次大型会议上谈论天主时的一些言论，说到我在关于圣神方面加上了那个引起喧嚷的措辞（我们要将灯藏在斗底下到何时呢？），\Quote{但另一个人却隐晦地暗示，只是稍微提及教义，并不公开说出真理；他用策略多于虔诚来灌满人们的耳朵，以口才的力量隐藏其歧见。}\par

\Quote{这是，}我说，\Quote{因为我（生活在偏僻角落，不为人知，大多数人不知道我说什么，甚至几乎不知道我说话）可以无危险地探讨哲学；但他的话语分量更重，因为他本人及其教会都更知名。他所说的一切都是公开的，围绕他的争战巨大，因为异端者试图抓住巴西略口中的每一个赤裸裸的词，将他赶出教会；因为他几乎是仅存的一星真理火花和生命力，周围的一切都被摧毁了；这样，邪恶就可能扎根于城中，并像从那教会的中心一样蔓延到全世界。因此，更好的做法是在真理上有所保留，我们自己像面对云彩一样对情况稍作让步，而不是因公开宣告而冒毁灭的危险。因为如果我们从其他导向这一结论的用语中承认圣神是天主（因为真理不在于声音而在于意义），那么我们就不会受到伤害；但如果真理因一个人的被驱逐而被赶走，那么教会就会受到极大的伤害。}在场的同伴不接受我的经济之策，认为它过时且嘲弄他们；他们大声反对我，说我实行它更多是出于怯懦而非理性。他们说，\Quote{用真理来保护我们自己的人，远比用你们所谓的\InnerQuote{经济}来削弱他们同时无法赢得外人要好得多。}这会是冗长且或许不必要的，详细告诉你我所说、所听的一切，以及我如何被对手惹恼（也许过度且违背我平素的性情）。但总之，我以同样的方式送走了他们。而您，神圣而尊贵的领袖，请教导我，在阐述圣神的天主性上我应走到多远；我应使用什么言辞，以及如何保留；好让我能武装起来对抗对手。因为如果连我，比任何人都更了解您和您的意见，且多次就此给予并接受保证，仍需要被教导此事真理，那么我将是所有人中最无知和可怜的。\par

% structure: p0072 source=h2 target=subsection reason=relative-heading-rank; base=h2

\subsection{书信五十九　致巴西略}

（对巴西略对上一封信有些愤怒的回应。）\par

这件事是任何一个更明智的人都会预见到的；但我这个非常单纯愚笨的人，在给您写信时并未担心。我的信使您忧伤；但在我看来，既不正确也不公正，而是完全不合理。您没有承认您受伤，也没有隐瞒，或者如果您隐瞒了，那也是非常巧妙，像一个面具，用尊敬的表面掩饰您的恼怒。至于我自己，如果我行事诡诈或恶意，我不只被您的恼怒惩罚，更被真理本身惩罚；但如果我是以单纯和惯常的善意行事，我将把责任归于我自己的罪过，而非您的脾气。但更好的做法是纠正此事，而不是对向您进谏的人发怒。但您必须处理自己的事务，因为您完全有能力给他人同样的劝告。您可以将我看作非常乐意，若天主愿意，前来与您会合，加入您的争战，并贡献我所能的一切。因为谁会退缩？谁不会在您领导下、在您身旁，为真理说话和争战时，反而鼓起勇气呢？\par

% structure: p0075 source=h2 target=subsection reason=relative-heading-rank; base=h2

\subsection{书信六十　致巴西略}

（额我略因母亲病重未能履行书信五十九末尾的承诺，故写信解释并致歉。）\par

完成您的吩咐部分取决于我；但部分，我敢说主要取决于您的尊威。取决我的是善意和热忱，因为我从未回避与您会面，而是一直寻求机会，此刻甚至更加渴望如此。取决您圣德的，是我的事务得以理顺。因为我正陪伴我的母亲夫人，她已患病多时。若我能离开她而无危险，您尽可确信我不会剥夺自己前往您那里的乐趣。因此，请以您的祈祷帮助我，求她康复，并祝我旅途顺利。\par

\SourceRecord{Translated by Charles Gordon Browne and James Edward Swallow.；Nicene and Post-Nicene Fathers, Second Series；Vol. 7.；Edited by Philip Schaff and Henry Wace.；Buffalo, NY: Christian Literature Publishing Co.,；1894.；<http://www.newadvent.org/fathers/3103b.htm>.}

% structure: p0001 source=h1 target=section reason=page-title

\section{书信（第三部分）}

\Emph{杂类书信}。\par

% structure: p0003 source=h2 target=subsection reason=relative-heading-rank; base=h2

\subsection{一、致其弟\Person{凯撒里}书}

% structure: p0004 source=h3 target=subsubsection reason=relative-heading-rank; base=h2

\subsubsection{第七封书信}

（\Person{君士坦提乌斯}皇帝去世后，无可争议的继承权落到了他的表弟、背教者\Person{尤利安}手中。\Person{尤利安}随即开始运用帝国的全部权力来阻挠基督信仰（尽管并未完全迫害），并恢复古异教的至高地位。他最初采取的行动之一，便是罢免所有前任皇帝手下担任高职的人。然而，\Person{额我略}的弟弟\Person{凯撒里}却是例外；\Person{尤利安}或许曾在雅典认识并器重他，所以竭尽全力将他留在宫廷，并拉拢他。这令\Person{凯撒里}以及他的其他友人——包括\Person{额我略}——深感忧虑。他们预见到\Person{尤利安}将竭力动摇这位年轻人的信德，而且不能确信他会有勇气抵抗这样的攻击。\Person{额我略}在烦恼中写了下面这封信给他。不久之后，预料中的试探果然来了。\Person{圣凯撒里}勇敢地对抗皇帝，在声明自己坚定不移地持守信德之后，辞去了宫廷的职务，退隐到\Place{纳西盎}。）\par

我为你感到羞愧已经够了；我忧虑自不必说，因你是我最了解的人。但且不提我个人的感受，也不提那关于你的传言带给我的痛苦（若容我说，还有恐惧），我本希望，若可能的话，你能听到别人——不论是亲戚还是外人，只要认识我们（当然我指的是基督内的弟兄）的人——关于你我的谈论；而且不只是一两个人，而是人人都依次如此说；因为人总是对陌生人比对自己的亲人更爱说道理。他们中间如此这般的话语已经成为一种练习：\Quote{如今主教的儿子从了军；如今他贪图外在的权势和名声；如今他做了金钱的奴隶，正当众人重新燃起大火、为生命奔跑的时候；而他却认为唯一的荣耀、安全与财富，就是在时代中高尚地站立，并使自己尽可能远离一切可憎与玷污。}那么，主教又怎能劝勉别人不要随波逐流，或者不要与偶像混同呢？他又怎能责备那些在其他方面犯罪的人，因为自己的家室剥夺了他直言的权利？我们每天都听到这些话，甚至更严厉的话，有些说话者或许是出于友谊，有些则带着敌意。你认为我们是什么感受？我们这些自称事奉天主、并以盼望未来为唯一善事的人，听到这样的话，心中作何状态？我们可敬的父亲因所听到的一切而极为痛苦，甚至厌世。我尽己所能安慰他、扶持他，为你的心志作担保，并向他保证你必不会这样继续使我们忧伤。但是，若我们亲爱的母亲听到了关于你的事（迄今我们靠各种办法瞒着她），我想她会完全无法安慰；因为她是妇人，意志软弱，且因其极大的虔诚，在这样的事上无法控制情感。你若对自己、对我们还有丝毫顾念，就寻求更美好、更稳妥的道路吧。我们的资产足以维持独立生活，至少对于一个有节制欲望、不贪得无厌的人而言是如此。而且，若我们错过了这个机会，我看不出我们还要等待什么时机让你安顿下来。但若你执意固执己见，认为一切都微不足道，只求满足自己的愿望，那我就不想再说什么使你不快的话了；只是我要预言并郑重声明：两者必居其一——要么你仍然作一个真正的基督徒，但将被列于最低微的人中，处于与你自己和你的希望不相称的地位；要么你追逐荣耀，却在更重要的事上伤害自己，即使不在火中，也必染上烟尘。\par

% structure: p0007 source=h3 target=subsubsection reason=relative-heading-rank; base=h2

\subsubsection{第十四封与第二十三封书信}

（在\Person{瓦伦斯}皇帝统治下，\Person{凯撒里}重返公共生活，被任命为\Place{比提尼亚}财务官。他在此任上时，他的兄弟写信给他，为两位表兄弟求情：一位是\Person{欧拉里}，后来继\Person{额我略}任\Place{纳西盎}主教，与\Person{额我略}关系亲密；另一位是\Person{安斐洛希}，他因合伙人的欺诈而在\Place{君士坦丁堡}遇到金钱纠纷，\Person{额我略}为他请求援助和建议。不过，有人认为这封信并非写给其兄弟（当时他可能已在\Place{君士坦丁堡}），而是写给宫廷中某位高官。\Person{安斐洛希}不久后隐退，到公元347年时已是重要的\Place{依科尼雍}教区主教。\Person{额我略}写给他的信件将在本卷后面部分给出。）\par

请向你和我施行一件善事，这种善事你不常有機會去做，因為這樣的機會不常出現。承担起对於我可爱的表兄弟的最正义的保护，他们为了一份他们购买作为退休安养、且能提供他们一些生活资料的产业而烦扰过多；但在完成购买之后，他们遭遇了许多麻烦，部分是由于发现卖主不诚实，部分是由于被邻居抢夺和抢劫，以至于他们若能以买入的价格——\Emph{再加上}他们额外花费的不小一笔钱——脱手这份产业，反而是有利的。因此，如果你愿意在审查契約之后，看如何最稳妥地办理，将此事转给自己，那么这对他们和我都是最可取的；但如果你不愿意，次佳的办法就是你去抵制那人的多管闲事和不诚实，使他不能凭借他们缺乏商业习惯而占到任何便宜，无论是如果他们保留产业时伤害他们，或是如果他们放弃产业时使他们遭受损失。我实在惭愧就如此之事写信给你。尽管如此，由于我们無論是出於他們與我們的親戚關係，還是出於他們的職業（因為有誰不願照顧这样的人呢？又有什麼比不願意施予這樣的恩惠更令人羞恥呢？），你都應當為你自己，或为我，或为他们自己，或为所有这些原因加在一起，无论如何向他们施行这件善事。\par

% structure: p0010 source=h3 target=subsubsection reason=relative-heading-rank; base=h2

\subsubsection{书信二十三}

若我向你要求一件大恩，不要惊讶；因为我正在向一位伟人要求，而请求必须以所请求之人的尺度来衡量；因为向一个小人要求大事，或向一个伟人要求小事，都是荒谬的，前者不合时宜，后者则卑劣。因此，我亲手向你推荐我最宝贵的儿子\Person{安斐洛邱}，他是一位以其君子风度而如此著名（甚至超越他的年龄）的人，以至于我自己，虽是一位老人、一位司铎、又是你的朋友，也完全愿意被如此尊重。他若被一个人的假友谊所欺骗，並且没有怀疑那骗局，又有什么奇怪呢？因为他自己不是恶棍，所以不怀疑恶行，反而以为需要的是言语的纠正而非品格的纠正，因此与他合伙经营。在正直的人眼中，这又能歸咎于他什么过错呢？所以不要让恶行胜过美德；也不要羞辱我的白发，而要尊重我的见证，并将你的恩惠加于我的祝福之上，这些祝福在我們所侍立的、天主面前或许有些价值。\par

% structure: p0012 source=h3 target=subsubsection reason=relative-heading-rank; base=h2

\subsubsection{书信二十}

（公元368年，比提尼亚的尼西亚城几乎被一场可怕的地震完全摧毁。\Person{凯撒留}失去了他的房屋，而他本人的逃脱几乎是奇迹。\Person{额我略}写信（如同\Person{巴西略}也做了）祝贺他的获救，並趁机敦促他從世俗事务中退隐。\Person{凯撒留}很快就决定遵循这一建议，并正采取措施实施这一决定时，于公元369年初突然去世，年仅40岁。他将所有巨额财产留给穷人，但这些财产一度落入策划者之手，作为他兄弟遗嘱执行人的\Person{额我略}，费了很大力气才为原本的目的收回这笔财产。（参见就此事写给君士坦丁堡总督\Person{索弗洛尼乌斯}的信。）他被安葬在纳西盎的殉道者教堂的一个墓穴中，那是他父母为自己预备的。\Person{额我略}发表了葬礼讲道，收录在本卷前部。这四封是已知的兄弟间仅有的书信。）\par

即使是惊恐对智者也不是没有用处；或者，我应该说，它们非常宝贵且有益。因为，尽管我们祈祷它们不发生，但当它们发生时，它们教诲我们。因为受折磨的灵魂，正如\Person{伯多禄}在某处精妙地说，亲近天主；而每一个逃脱危险的人会更亲近于那保存他的主。因此，让我们不要因我们分担了灾祸而烦恼，而要为得救而感恩。并且，让我们不要在危险时显示自己是一样，在危险过后又是另一样，而要下定决心，无论在家在外，无论私人生活或公職（因为我必须说这一点，不可省略），都要跟随那保存我们的主，依附于祂的身旁，少思想世上微小的事；让我们为后来的人提供一个故事，为我们的荣耀和灵魂的益处，同时也为所有人提供一个非常有用的教训，即危险比安全更好，不幸比成功更可取，至少如果在恐惧之前我们属于世界，但在恐惧之后我们属于天主。也许我对你来说有些厌烦，因为我如此频繁地就同一主题写信给你，你会认为我的信不是劝勉而是炫耀，所以到此为止。你会知道，我尤其渴望能与你在一起，分享你获救的喜悦，并在稍后谈论这些事。但既然这不可能，我希望尽快在这里接待你，共同庆祝我们的感恩。\par

% structure: p0015 source=h2 target=subsection reason=relative-heading-rank; base=h2

\subsection{致\Person{尼撒的圣额我略}}

（\Person{尼撒的额我略}主教是\Person{大巴西略}的弟弟。早年被读经职，但他厌倦了圣召，成为修辞学教授。这在教会中引起了丑闻，并使他的朋友大为悲伤。\Person{纳西盎的额我略}写了以下劝诫信，这并非没有效果，因为不久之后他就放弃了世俗的职份，退隐到他的兄弟\Person{巴西略}在庞图斯建立的隐修院。他在那里花了几年时间研究圣经和最好的注释家。）\par

% structure: p0017 source=h3 target=subsubsection reason=relative-heading-rank; base=h2

\subsubsection{书信一}

在我性格中有一个优点，我要在许多点中夸耀这一点：我对一个坏计划同样会为自己和朋友感到烦恼。既然所有按照天主生活、遵行同一福音的人都是朋友和亲属，那么为什么你不应该从我这里直白地听到所有人悄悄说的话呢？他们不赞同你那不光彩的荣耀（借用你自家行当的一个说法），不赞同你逐渐堕落到更低级的生活，也不赞同你的野心——按照欧里庇得斯的说法，这是最坏的恶魔。最明智的人啊，你遭遇了什么，你为何自我谴责，竟抛弃了你曾经向民众诵读的那些神圣而令人喜悦的书籍（不要羞于听到这个），或者把它们像冬天人们对待舵和锄头一样挂在烟囱上，而投身于咸涩苦涩的书籍，宁愿被称为修辞学教授而非基督信仰教授？我感谢天主，宁可做后者而非前者。我亲爱的朋友，不要再这样下去了，尽管为时已晚，但要恢复清醒，重新回归你自己，向信友、向天主、向祂的祭台和圣事——你曾从这些中退却——做出你的道歉。不要以傲慢的修辞风格对我说：怎么，我教修辞学时难道不是基督徒吗？我和孩子们在一起时难道不是信徒吗？也许你还会呼求天主作证。不，我的朋友，即使我部分地同意你，你也没有达到你应当达到的彻底程度。你那当前职业给别人——那些天生倾向于邪恶的人——带来的冒犯，以及给他们机会去思考和说出你最坏的一面，又该如何？我承认是虚假的，但何必要这样呢？因为人不是只为自己活，也为了邻居；仅说服自己是不够的，还必须说服别人。如果你在公开场合练习拳击，或在剧院里击打和挨打，或羞愧地扭动身体，你会说自己拥有节制的灵魂吗？这样的论据不适合明智的人；接受它是轻浮的。如果你做出改变，我现在就会欢喜，正如一位毕达哥拉斯派的哲学家在哀叹朋友堕落时所说；但他写道：如果不，你就对我而言是死了。但为了你，我还不会这样说。\Quote{曾是朋友，却成了敌人，但仍是朋友，}正如悲剧所说。如果既不能自己看见正确的事——这是最高的方法，又不愿听从他人的建议——这是次好的，那我会感到悲伤（说得温和些）。至此是我的劝告。请原谅我对你的友谊使我悲伤，并使我为你、为整个司铎团，我还可以补充说，为所有基督徒而激动。如果我可以与你一同祈祷或为你祈祷，愿那使死人复生的天主援助你的软弱。\par

% structure: p0019 source=h3 target=subsubsection reason=relative-heading-rank; base=h2

\subsubsection{第七十二书}

（当圣格列高利被祝圣为尼撒主教时，帝国皇位由瓦伦斯占据，他是一位热心的亚略派信徒，一心要摧毁尼西亚信仰。他为此任命了一位名叫德摩斯梯尼的前御厨房文书为庞迪克民政教区的代理总督。与巴西略的旧怨使此人对格列高利不友好，在用小方式迫害了他一段时间后，公元275年，他召集了一次主教会议，调查其祝圣仪式中一些违规的指控，并因一些管理教会资金不当的轻浮罪名审判格列高利。格列高利因患胸膜炎未能出席在安基拉举行的这次会议；另一次会议被召集在尼撒本地举行。但格列高利拒绝出席，并被以藐视法庭罪撤职。于是瓦伦斯放逐了他，他似乎陷入了非常低落的情绪，几乎因异端派别的明显得势而沮丧。随后的三封信揭示了他此时的一些心境。这些信是回复他现已失传的信件，目的是安慰他的困境，鼓励他重新振作，盼望好日子到来。这种更加乐观的语气被后来的事件证明是合理的：公元378年瓦伦斯去世后，被放逐的主教们被格拉提安恢复，格列高利重登主教宝座，令其教区的信友大为欢喜。）\par

不要让你的烦恼过度困扰你。因为我们越不为事情悲伤，事情就越不令人悲伤。异端者们已经解冻，借着春天壮起胆来，从他们的洞里爬出来，如你所写，这并不奇怪。我知道，他们会嘶叫一阵，然后又会躲藏起来，被真理和时势所征服，而且我们越是将整个事情交托给天主，他们就越是如此。\par

% structure: p0022 source=h3 target=subsubsection reason=relative-heading-rank; base=h2

\subsubsection{第七十三书}

关于你信中提到的主题，这是我的感受。我不因被忽视而生气，但我因被尊敬而欢喜。前者是我应得的，后者是你尊重的证明。请为我祈祷。原谅这封短信，因为无论如何，尽管它很短，却比沉默要长。\par

% structure: p0024 source=h3 target=subsubsection reason=relative-heading-rank; base=h2

\subsubsection{第七十四书}

虽然我在家，但我的爱与你一同在外流亡，因为亲情使我们万物共有。信赖天主的仁慈和你的祈祷，我怀着极大的希望，一切都会按照你的心意成就，暴风将变成和风，天主将因你的正统信仰赐予你这奖赏，使你胜过你的对手。我最渴望早日见到你，与你共度美好时光，正如我所祈祷的。但如果你因事务缠身而延迟，至少用一封信使我欢欣，并不要不屑于告诉我你的一切情况，并像你习惯的那样为我祈祷。愿天主赐予你在一切情况下健康与喜乐——你是整个教会的共同支柱。\par

% structure: p0026 source=h3 target=subsubsection reason=relative-heading-rank; base=h2

\subsubsection{第七十六书}

（大巴西略于公元379年1月1日去世。纳齐安的格列高利因重病未能参加他的葬礼，因此写信给尼撒的格列高利如下。）\par

此事亦为我这悲伤的人生所留，即听闻巴西略的死讯，以及那圣善灵魂的离去——它离开我们，为能与主同在，这正是他毕生所预备的。在我不得不承受的一切其他损失中，这是最大的：因我仍受身体疾病的折磨，且处于极大危险中，无法亲吻那圣洁的尘土，也无法与你们同在，享受真正哲学的慰藉，并安慰我们共同的朋友。然而，眼见教会如此荒凉，失去如此荣耀，被夺去如此冠冕，这是任何稍有感情的人，都无法忍心观看或听闻的。但我想，你虽有许多朋友，并将收到许多慰问之辞，但你的安慰更多来自你自己和对他的记忆，而非来自他人；因你们二位是众人哲学的表率，一种属灵的准绳，既体现顺境中的纪律，也体现逆境中的坚忍；因为哲学以节制承载顺境，以尊严承载逆境。这是我对阁下要说的话。至于我写此信的人，除了你的陪伴和交谈，还有什么时间或言语能安慰我呢？我们蒙福的那位已将你留给我，代替一切，使我在你身上如同在明亮闪耀的镜子中看见他的品格，从而认为自己仍拥有他！\par

% structure: p0029 source=h3 target=subsubsection reason=relative-heading-rank; base=h2

\subsubsection{Ep. LXXXI.}

你因旅途劳顿而苦恼，认为自己漂泊不定，如同被水流带动的一根树枝。但我亲爱的朋友，你万万不可如此感觉。树枝的漂流是不由自主的，而你的行程是由天主所安排，你的稳固在于为他人行善，即使你并不固定于一处；除非有人应当责备太阳，因为它周游世界，散布光芒，并赐予它所照耀的一切事物生命；或者，有人在赞美恒星的同时，也应谴责行星，它们的游走本身却是和谐的。\par

% structure: p0031 source=h3 target=subsubsection reason=relative-heading-rank; base=h2

\subsubsection{Ep. CLXXXII.}

（额我略在辞去君士坦丁堡宗主教职位后，隐退至纳齐盎，并被说服承担当时空缺教区的管理，直至空缺被填补。该省主教们希望他永久保留此职，因此并不急于进行选举。但最终，他们屈服于额我略不断表达的意愿，选举了他的表亲尤拉利乌斯。然而不久，额我略的敌人散布谣言，说此次选举是违背他意愿的，意图不公正地将他驱逐出该教会的管理。以下信件即因这一诽谤而作。）\par

我有祸了，因我的寄居拖延了，而最大的不幸是，我们中间有战争与纷争，我们没有持守从圣父们所领受的和平。我深信你必藉着那扶持你和你所属之人的圣神的能力，恢复这和平。但恳请任何人，不要散布关于我和我诸位主教大人的虚假传言，仿佛他们违背我的意愿宣布了另一位主教来接替我的位置。然而，由于我身体虚弱，深陷困境，并害怕忽略教会的责任，我请求他们赐予我这恩惠——这并不违反教规，且对我是一种解脱——即为教会提供一位牧者。他已被赐予你们的祈祷，一位配得上你们虔敬的人，我现在将他交在你们手中，即至可敬的尤拉利乌斯，一位非常爱天主的主教，我愿在他的怀中死去。若有人认为在一位主教在世时另立主教是不妥当的，当知道在这件事上他无法抓住我们的把柄。因为众所周知，我被任命并非为纳齐盎，而是为萨西玛，尽管因尊敬我父亲，我曾以客人的身份在短时间内承担了管理之职。\par

% structure: p0034 source=h3 target=subsubsection reason=relative-heading-rank; base=h2

\subsubsection{Ep. CXCVII. 致悼其妹塞奥塞比亚之慰问信函}

（《传记辞典》中论尼撒的额我略一文的作者推测她是他的妻子，但除本信中称她为\Quote{一位司铎的真正伴侣}这一模糊表述外，并无其他证据；然而在同一封信中，她又明确被称为他的妹妹。一些作者想象她是纳齐盎的额我略本人的妻子，但没有证据表明他结过婚。她去世的日期不确定，但可能在公元381年之后。似乎\Quote{伴侣}一词可泛用于在同一工作中共享之人，因此本笃会编者认为塞奥塞比亚是尼撒教会的一位女执事。）\par

我本已匆忙启程前赴你处，且已行至\Place{尤斐弥亚斯}，却被你们为纪念圣殉道者所举行的庆节耽误了；一方面因为我身体欠佳，未能参与；另一方面，在如此不适宜的时节来临，可能会给你带来不便。我启程，部分是为了在阔别之后能见到你，部分是为了能钦佩你在你那位圣洁有福的姐妹离世时所展现的忍耐与哲学（因我已听闻此事），你好似一位善良而完美的天主的仆人，比任何人都更明了天主与人的事；你将那对他人而言极为沉重的事视为轻忽之事，即与这样一位灵魂共处，又送她离去，将她像打谷场上按时收积的禾捆一样储藏在安全的仓廪里，借用\WorkTitle{圣经}的话（\BibleRef{约 5:26}）；而且在她因生命短暂尝到人生喜乐、避开忧伤之后，在她尚未为你披麻戴孝之前，便受到你给予那配得如此之人的美好葬礼荣耀。我亦，信我，渴望离世，即使不及你那般（此言甚为夸张），也仅稍逊于你。然而面对天主长久有效的律法，它如今已带走了我的\Person{特奥塞维亚}（我之所以称她为我的，是因她度了虔敬的生活；因为精神的亲属关系优于肉体的），\Person{特奥塞维亚}，教会的荣耀，基督的装饰，我们这代的助佑，女性的希望；\Person{特奥塞维亚}，在众弟兄的所有美丽中至为美丽至为荣耀；\Person{特奥塞维亚}，真正神圣，真正是司铎的配偶，同样尊荣，堪当伟大的圣事；\Person{特奥塞维亚}，所有未来时代都要接纳她，她安息在不朽的柱石上，即所有现今认识她的人和所有后世之人的灵魂上。你勿惊奇我屡屡呼唤她的名字。因为我甚至因怀念这位有福者而喜乐。让这些，言简意赅，作为我致她的墓志铭，也是我对你的慰问之言，尽管你自己因在一切事上的哲学而能以这种方式安慰他人。我们相会（我极其渴望）因我前述的理由而受阻。但我们彼此祈祷，只要还在世上，直到共同的终末——我们正趋近的终末——追上我们。因此我们必须忍受一切，因为我们不会长久享有喜乐或遭受苦难了。\par

% structure: p0037 source=h2 target=subsection reason=relative-heading-rank; base=h2

\subsection{三、致\Person{欧瑟伯}，\Place{萨莫萨塔}主教}

% structure: p0038 source=h3 target=subsubsection reason=relative-heading-rank; base=h2

\subsubsection{第四十二书}

（本函催促他的朋友出席\Place{凯撒利亚}，以选举继任\Person{欧瑟伯}的都主教，已在本选集的第二部分给出。）\par

% structure: p0040 source=h3 target=subsubsection reason=relative-heading-rank; base=h2

\subsubsection{第四十四书}

（\Person{欧瑟伯}应上述呼吁动身前往\Place{凯撒利亚}后，\Person{老额我略}虽身体极其虚弱，仍决定亲自出席主教会议，以期通过他们共同的努力确保\Person{巴西略}的当选。\Person{小额我略}托其父带去以下信函，向他的朋友解释他本人未能前来的原因。日期约在公元三七九年九月。）\par

我从何处开始赞颂你，又以何名称呼你才是合宜的？教会的柱石和基础，或世上的光（借用宗徒的话），或基督信仰世界剩余部分的荣耀之冠，或天主的恩赐，或祖国的壁垒，或信德的旗帜，或真理的使者，或同时是这一切，甚至超越这一切？这些过分的赞颂，我将凭我们将要看到的事加以证明。有哪场雨曾如此及时地降在干涸的土地上？有哪股从岩石流出的水曾如此及时地供给旷野中的人？有哪样天使的食粮曾被人吃过？共同的主耶稣何曾如此及时地出现在将要溺毙的门徒面前，平息大海，拯救那行将灭亡的人，正如你在我们疲惫困苦、仿佛即将遭遇船难之际，向我们显现的那样？无需多言其他，你以何等勇气和喜乐充满了正统者之灵魂，你使多少人摆脱了绝望。\par

然而我们的母亲教会——我指的是\Place{凯撒利亚}——如今因见到你，真的在脱去她的寡妇衣袍，重新披上喜乐的外袍，并且当她领受一位配得上她自己、配得上她先前主教们、也配得上你双手的牧者时，她将更加光彩照人。因为你自己看到我们的事务状况如何，你的热忱、辛劳和敬虔的直言不讳成就了何等奇迹。年岁更新，疾病被战胜，卧床者跳跃，软弱者束上力量。由此我猜测，我们的事务也将如我们所愿地成就。你还有我的父亲，他代表他自己和我，要藉着当前为教会的奋斗为他整个一生和他可敬的晚年画上光荣的句号。我深信，他将因你的祈祷而坚固，青春更新地回到我身边；因为人必须凭信德勇敢地将一切托付于他们。然而，若他在此焦虑中结束生命，那达到如此结局亦非灾祸。恳求你，宽恕我，若我稍微向恶人的舌头让步，稍稍延迟前来拥抱你，并亲自完成我现在略过的应归于你的赞美。\par

% structure: p0044 source=h3 target=subsubsection reason=relative-heading-rank; base=h2

\subsubsection{第六十四书}

（三七四年，\Person{欧瑟伯}与其他东方正统主教被\Person{瓦伦斯}流放，他们的教座被亚略派侵入者占据。\Person{欧瑟伯}奉命退隐至\Place{色雷斯}，他的行程经过\Place{卡帕多西亚}，在那里他见到了\Person{巴西略}。但\Person{额我略}极为悲伤的是，他身体太不适，无法离家去会见他。于是他送去了以下信函。）\par

当尊驾经过我国时，我病重得甚至无法\Emph{向外看}。我悲伤，与其说是因为疾病（尽管它带来了对最坏后果的恐惧），不如说是由于未能与您的圣德与良善相见。我渴望见到您可敬的面容，就像一个人自然渴望从您那里得到属灵伤口的医治，并期待从您获得治愈一样。但当时，由于我的罪过，我错过了与您的会面；现在，藉着您的良善，我有可能为我的烦恼找到补救，因为如果您肯屈尊在您蒙悦纳的祈祷中纪念我，这将成为我在今世和来世从天主那里获得一切祝福的宝藏。因为像您这样的人，福音信仰的斗士，忍受了如此多的迫害，并因在患难中的忍耐为自己赢得了在至公义的天主面前的信心——这样的人若肯屈尊也在他的祈祷中作我的主保，我相信，这将为我赢得如同通过一位神圣殉道者所获得的同样多的力量。因此，请允许我恳求您，在所有我希望值得您纪念的事上，不间断地纪念您的\Person{额我略}。\par

% structure: p0047 source=h3 target=subsubsection reason=relative-heading-rank; base=h2

\subsubsection{Ep. LXV.}

（\Person{欧瑟伯}已回复前一封信，\Person{额我略}再次写信，藉着\Person{欧瑟伯}的一位门徒\Person{欧帕克西乌}得到与他朋友交流的机会，\Person{欧帕克西乌}在前往拜访他老师的途中经过\Place{卡帕多西亚}。此信有时归为\Person{巴西略}所作。）\par

我们可敬的弟兄\Person{欧帕克西乌}一向是我所亲爱的和真诚的朋友，但藉着他对你所怀的亲情，他显得更亲爱、更真诚，因为即使现在，他也匆匆赶到你尊前，如同用\Person{达味}的话说，像一只鹿奔向甜美的泉源，以解他那巨大且难以忍受的干渴——你患难中的忍耐。请屈尊作他和我的主保。\par

那些被允许接近你的人确实有福，更有福的是那能为他为基督所受的苦难和为真理所付的辛劳戴上冠冕的人，这冠冕是敬畏天主的少数人所获得的。因为你所显示的并非未经考验的德行，也并非只在风平浪静时你才正确航行并引导他人灵魂，而是在诱惑的艰难中你闪耀发光，比你的迫害者更伟大，并崇高地从你出生之地离去。他人拥有祖传的门槛——我们拥有天上的城；他人或许拥有我们的宝座——我们拥有基督。哦，何等有益的交换！我们舍弃何其少，而得何其多！我们经过了火与水，我相信我们也必将进入安息之所。因为天主不会永远舍弃我们，也不会将真信仰交于迫害，但按我们痛苦的众多，祂的安慰必使我们欢喜。至少我们相信并渴望如此。但请你，我恳求，为我们卑微之人祈祷。每当有机会，请毫不犹豫地以书信祝福我们，并以您的消息鼓舞我们，就如你刚恩赐所做的那样。\par

% structure: p0051 source=h3 target=subsubsection reason=relative-heading-rank; base=h2

\subsubsection{Ep. LXVI.}

（以下书信有时归为\Person{巴西略}所作，并存于他的作品及\Person{额我略}的作品中。然而，除一份手稿外，其余手稿均将其归于后者。）\par

你以写信和纪念我使我喜悦，更因在信中赐我祝福使我大大喜悦。但若我配得上你的苦难和你为基督并藉着基督的争战，我必配得上前去就你，拥抱你的虔诚，并以你在苦难中的忍耐为榜样。然而，我既然不配，因多受困苦和阻碍，我便做次好的事：我致书你的成全，求你莫厌倦纪念我。因为被认为配得上你的书信，于我不仅有益，且是可向许多人夸耀的事，也是荣誉，因我被如此德行高卓、与天主关系密切之人所看重，以致他能以言以行引领他人也进入与天主的关系。\par

% structure: p0054 source=h2 target=subsection reason=relative-heading-rank; base=h2

\subsection{四、致\Place{君士坦丁堡}总督\Person{索弗洛尼乌斯}}

（\Person{索弗洛尼乌斯}，\Place{卡帕多西亚的凯撒勒亚}人，是\Person{额我略}和\Person{巴西略}的早期朋友和同学。他进入文官系统，很快升至高位。公元365年，他因及时向\Person{瓦伦斯}皇帝通报\Person{普罗科皮乌斯}的篡位企图，被任命为\Place{君士坦丁堡}总督。他主要为\Person{额我略}和\Person{巴西略}的信件所提及，这些信件请求他为各种人物施以援手。第21封信写于公元369年，向\Person{索弗洛尼乌斯}推荐\Person{尼科布禄}，\Person{额我略}的侄婿，\Person{阿利比亚娜}的丈夫，他姊姊\Person{果尔果尼亚}的女儿。这位\Person{尼科布禄}极富才干，但极不喜公职。\Person{额我略}不断写信给一位又一位高官，以求为他免除强加于他的任命。）\par

% structure: p0056 source=h3 target=subsubsection reason=relative-heading-rank; base=h2

\subsubsection{Ep. XXI.}

黄金在不同时期被改变和转化，塑造为多种装饰，被艺术用于多种用途；但它仍是它自己——黄金；变化的是形式而非本质。同样，我相信您的仁慈对您的朋友将保持不变，尽管您不断攀升高位，我冒昧向您呈上此请求，因为我不仅尊崇您的高位，更信赖您的善心。我恳求您对我极其可敬的儿子\Person{尼科布禄}施以恩惠，他在各方面都与我联结——因亲属关系，因亲密交往，更重要的是因性情。在什么事上，到什么程度？在他可能请求您帮助的任何事上，并限于您认为符合您的宽宏大量的程度。至于我，将以我所有的最好的来回报您：我有言辞的能力，能够宣扬您的良善，即使不能完全按您的价值，至少尽我所能。\par

% structure: p0058 source=h3 target=subsubsection reason=relative-heading-rank; base=h2

\subsubsection{Ep. XXII.}

（此信为\Person{安斐洛基}而写，写于同一时间，并因与致\Person{凯撒留斯}书信中置于第二封者相同的困境而作。）\par

正如我们通过外观辨识金银与石头，同样地，我们也能以此区分好人与坏人，无需长久的考验。因为我不需要很多言语，就我最高贵的儿子\Person{安非洛基}向您的高贵求情。我更宁愿期望发生某种奇怪而难以置信的事，也不愿他在金钱问题上做出任何不名誉的事，或起这样的念头；他作为君子，且比他的年龄更智慧，享有如此普遍的声誉。但他必须承受什么？没有什么能逃脱嫉妒，因为一些指责之言也触及了他，一个因单纯而非心性堕落而陷入犯罪指控的人。但请不要容许你在他烦恼和困苦中忽视他。不，我恳求你神圣而伟大的心，但请荣耀你的祖国，扶助他的德行，并尊重我——我藉着你并透过你获得了荣耀；成为这个人的一切，将意愿加于能力，因为我知道没有什么能与阁下的权力相提并论。\par

% structure: p0061 source=h3 target=subsubsection reason=relative-heading-rank; base=h2

\subsubsection{书简二十九}

（同年。此处\Person{凯撒里乌}曾将其全部财产遗赠给穷人；但他的房子被仆人洗劫，他的朋友们只找到一笔相对较小的金额。除此之外，不久后一些人自称是其遗产的债权人，尽管无法证明，他们的要求仍被支付。然后其他人不断出现，直到最后家人拒绝再承认任何要求。接着有人威胁要诉讼。\Person{额我略}极度厌恶这一切，并且担心诉讼可能引起的丑闻，于是致函总督如下。）\par

您看我的处境如何，人类事务的循环如何运转，时而一些人兴旺，时而另一些人衰败，如俗话所说，兴衰无常，总是变化更迭，以至于人宁可相信微风，或写在水上的文字，也不信人的顺遂。这是何故？我想，这是为了藉着默观这一切的不确定与变幻，我们学会更多地投奔天主与未来，对阴影和梦境稍加留意。但什么引起了这番话？因为我并非无故如此哲学思考，也非徒然夸耀？\par

\Person{凯撒里乌}曾是你并非最不著名的朋友之一；确实，除非兄弟之爱欺骗了我，他是你最杰出的朋友之一，因为他学识渊博，品行超群，且以其众多朋友而闻名；在这些朋友中，正如他一直认为并使我确信的，阁下位居首位。这些都是旧事，你在纪念他时自会锦上添花；因为人之常情是为逝者增添赞美。但现在（免得你无泪地略过此事，或使你流泪有益），他死了，无友，孤单，可怜，仅得到一点没药（如果有的话），以及最后的小覆盖物，他能得到这点怜悯已是幸事。但他的仇敌，如我所闻，已扑向他的遗产，从四面八方猛烈掠夺，或即将如此。残酷啊！野蛮啊！无人阻止；即使最善良的朋友也只是以法律作为最大的恩惠。简言之，我曾是幸福之人，如今却成了一出戏剧。请不认为此事可容忍，而是以同情和与我同愤来帮助我，为死去的\Person{凯撒里乌}伸张正义。是的，以友谊本身的名义；是的，凭你最珍视的一切；凭你的希望（愿你通过向逝者展示信实和真诚而使其稳固），我恳求你向活着的人行此善事，使他们怀有希望。你以为我因金钱而悲伤？如果\Person{凯撒里乌}独自一人——他曾以为有这么多朋友——结果却一个都没有，那将是我更难以忍受的耻辱。这就是我的请求，以及产生此请求的原因，因为我的事对你并非完全无关紧要。你将如何协助我，用什么方式，以及如何，事情本身会提示，你的智慧会考虑。\par

% structure: p0065 source=h3 target=subsubsection reason=relative-heading-rank; base=h2

\subsubsection{书简三十七}

（一封推荐信，为修辞学家\Person{欧多克修斯}，\Person{额我略}对他怀有深厚的敬意。）\par

孝顺母亲是宗教义务。现在，不同的人有不同的母亲；但所有人的共同母亲是我们的祖国。你以你整个生活的光辉荣耀了这位母亲；现在你将再次荣耀她，藉着为我获得我所恳求的。我的请求是什么？你肯定认识修辞学家\Person{欧多克修斯}，他是她儿子中最博学的。简言之，他的儿子，另一个\Person{欧多克修斯}，无论在生活还是学识上，现在通过我来接近你。为了使你获得更好的名声，请在他请求你协助的事务上帮助他。因为如果你，我们祖国的普遍保护者，已经对许多人行了善，我还要加上，将继续行善，却不优先尊敬那位在学识和口才上最卓越的人，那将是一种耻辱，你应该尊敬他，即使没有其他理由，只因为他用口才赞美你的善德。\par

% structure: p0068 source=h3 target=subsubsection reason=relative-heading-rank; base=h2

\subsubsection{书简三十九}

（大约同一日期。推荐一位\Person{阿马佐尼}，\Person{额我略}非常尊重他的学识。）\par

我衷心祝愿我所有的朋友安好。当我说朋友时，我指的是那些与我在德性上相连的可敬而善良的人，倘若我本人尚有任何德性可言的话。因此，在当下我寻求如何善待我杰出的兄弟\Person{亚马佐尼乌斯}时（因为最近我们之间有过一些交往，我对他非常满意），我想我可以以一还百——那便是回报他以你的友谊和保护。因为在短时间内，他便展示出广博学识的明证，既有我过去曾热心追求的那一类（当时我目光短浅），也有我自能默观德性巅峰以来按其时宜所热心追求的那一类。我是否反过来在他眼中在德性上有什么价值，那是他的事。无论如何，我向他展示了我拥有的最美好之物，即我的朋友们，视他如朋友。在这些朋友中，我认为你是首要且最真诚的一位，我愿你对他也是如此——正如我们共同的祖国所要求，也如我的愿望和承诺所恳求；因为我曾应许他你的庇护，以回报他全部的善意。\par

% structure: p0071 source=h3 target=subsubsection reason=relative-heading-rank; base=h2

\subsubsection{Ep. XCIII.}

（写于额我略辞去总主教职务后不久。）\par

我们的退隐、闲暇与宁静对我而言颇为惬意；但它们使我与你的友谊和交往隔绝，这一点却并非有益，反倒不然。别人享有阁下的完德，而对我而言，若能拥有哪怕只是信件中带来的一丝交谈的影子，也实为一大幸事。我还能再见到你吗？我还能再拥抱这位我引以为荣的人吗？这些会恩赐给我余生的时日吗？若是如此，一切感谢天主；若非如此，我生命最好的部分已经结束。请记得你的朋友额我略，并为他祈祷。\par

% structure: p0074 source=h3 target=subsubsection reason=relative-heading-rank; base=h2

\subsubsection{Ep. CXXXV.}

（大约在公元382年中，狄奥多西在圣达玛苏的推荐下，召集了一次东方主教的新会议在君士坦丁堡举行，试图治愈因安提约基雅选举弗拉维安而加剧的分裂。额我略一听说这次会议召开，便写信给他在朝廷的几位有影响力的朋友，恳请他们为促进和平尽最大努力。）\par

我正悠闲地哲学思辨。这就是我的敌人给我带来的伤害，我巴不得他们再多做些这样的事，好让我更加视他们为恩人。因为常常发生这样的事：受委屈的人得到益处，而我们要善待的人却受到伤害。我的情况就是如此。但如果我不能让每个人都相信这一点，我则迫切希望至少你——我最愿意向其叙述我境况的人——能知道，更确切地说，我确信你已经知道，并能说服那些不知道的人。然而，我恳求你，现在无论如何（如果你以前没有做过）要尽力使那些不幸分裂的世人在声音和心意上合而为一；尤其是如果你觉察到（正如我所观察到的）他们分裂并非由于信仰，而是由于琐碎的私利。若能做到这一点，你将获得赏报；而我的退隐也会少一些悲伤，如果我能看到自己并非徒然紧握不放手，而是像一位约纳，甘愿投入海中，好使风暴止息，水手得救。但如果他们依然像以往一样颠簸不定，我至少已尽己所能。\par

% structure: p0077 source=h2 target=subsection reason=relative-heading-rank; base=h2

\subsection{五、致小安菲洛奇乌斯}

% structure: p0078 source=h3 target=subsubsection reason=relative-heading-rank; base=h2

\subsubsection{Ep. IX.}

（君士坦丁和君士坦提乌斯曾赐予所有圣职人员免除军事税的特权。然而，这一特权被尤利安废除，后由瓦伦提尼安和瓦伦斯恢复；但税务官常不顾豁免仍企图向他们征税。纳齐盎的税务官试图对一位名叫欧塔利乌斯的执事征税，额我略便代他写信给当时担任该省主要长官的安菲洛奇乌斯。此信日期约为公元372年，即额我略被祝圣为司铎之年。关于这位安菲洛奇乌斯的更多细节，见致凯撒留第二、第三信的导言（书信22、23）。）\par

如同\Person{品达}所说，以金柱支撑建造良好的厅室，请从起初在当前的忧虑中向我们展示正确的一面，好为你自己建造一座显赫的宫殿，并在其中以美名显现。但你将如何做到？通过尊崇天主和天主的事物，在你眼中没有比祂更伟大的。但如何、并以何种行动尊崇祂？通过这一行动：保护天主的仆役和祭台的侍奉者。其中之一是我们的同道执事欧塔利乌斯，不知为何，行政官员试图在他晋升更高品级后向他征收金币。请勿允许此事。向这位执事、向全体圣职人员、尤其向我（你所关心的人）伸出援手；否则他将独自一人承受不公正的对待，被剥夺时代的恩惠和皇帝赐予圣职人员的特权，甚至可能因我的软弱而受辱并被罚款。即便他人不怀好意，你来阻止此事也是合适的。\par

% structure: p0081 source=h3 target=subsubsection reason=relative-heading-rank; base=h2

\subsubsection{Ep. XIII.}

（见致索夫罗尼乌斯的第一封信。此处提及的麻烦性质不详。有几封给不同的人的信涉及他的苦恼和困境，其中许多源于他不愿承担任何公职的责任。他英年早逝，留下遗孀阿利皮安娜带着一大家人在极其窘迫的境况中度日。她的苦恼和子女的教育是额我略非常关心的事，他频繁的信件将见下文。）\par

我赞同\Person{泰奥格尼斯}的陈述，他虽然不赞美仅限于杯酒和欢乐的友谊，但赞美延及行动的友谊，用以下话语：\par

\Quote{满杯酒旁，人人皆友；\newline{}艰难临头，挚友难求。}\par

然而，我们彼此并未共饮一杯，甚至也不常会面（虽然我们本应十分谨慎地这样做，既为我们自己，也为我们从父辈继承的友谊），但我们确实请求那表现为行动的好意。一场斗争迫在眉睫，而且非常严重。我的儿子\\Person{Nicobulus}遇到了意外的麻烦，来自最不可能发生麻烦的地方。因此我恳求你尽快来帮助我们，既参与审理此案，若你发现我们受到不公正对待，也为我们辩护。但如果你不能来，无论如何也不要让自己被对方预先聘去，或者为了微小的利益出卖那据众人见证始终属于你的自由。\par

% structure: p0086 source=h3 target=subsubsection reason=relative-heading-rank; base=h2

\subsubsection{《书信》第二十五封}

\\EditorialNote{（\\Person{Amphilochius}在先前书信所提及的指控中被宣告无罪；但这次控告在他心中产生的结果是，他辞去了职务，退隐到一处名为\\Place{Ozizala}的隐修所，离\\Place{Nazianzus}不远，在那里他将劳动时间用于种植蔬菜。以下四封信没有特别的重要性，仅作为\\Person{Gregory}能与其密友使用的轻松风格的样本。）}\par

我没有向你求饼，正如我也不会向\\Place{Ostracine}的居民求水。但如果我向一个\\Place{Ozizala}的人求蔬菜，那并不奇怪，也不会对友谊造成过重的负担；因为你那里有很多蔬菜，而我们这里却很匮乏。所以我求你送一些蔬菜给我，要多多的，最好的品质，或者尽可能多（因为即使一点点东西对穷人来说也是巨大的）；因为我将要接待伟大的\\Person{Basil}，而你曾体验过他，既丰盛又富有哲学，你不会愿意看到他饥饿而烦躁。\par

% structure: p0089 source=h3 target=subsubsection reason=relative-heading-rank; base=h2

\subsubsection{《书信》第二十六封}

你送我的蔬菜真是太少了！它们一定是黄金蔬菜！然而你的全部财富在于果园、河流、树林和花园，你的土地出产蔬菜如同其他土地出产黄金，而且\par

你住在草木繁茂的叶丛之中。\par

但谷物对你们来说是神话般的幸福，你们的面包是天使的面包，正如俗话所说，如此受欢迎，而又如此不可靠。那么，要么你大方地送蔬菜给我，要么——我不威胁你别的，但我不送任何谷物给你，看看是不是真有\\Quote{蝗虫靠露水为生}这句话！\par

% structure: p0093 source=h3 target=subsubsection reason=relative-heading-rank; base=h2

\subsubsection{《书信》第二十七封}

你在开玩笑；但我知道，一个\\Place{Ozizala}人，当他最辛勤劳作时，却饿死的危险。只给与他们这样的赞美：即使他们饿死，也气味芬芳，并有华丽的葬礼。怎么会这样？因为他们被各种鲜花覆盖。\par

% structure: p0095 source=h3 target=subsubsection reason=relative-heading-rank; base=h2

\subsubsection{《书信》第二十八封}

在访问与\\Place{Pamphylia}接壤的山城时，我在山中捞起了一只海中的\\Person{Glaucus}；我不是用亚麻网把鱼从深水里拖出来，而是用朋友的爱捕捉了我的猎物。既然我已经教会我的\\Person{Glaucus}在陆地上旅行，我派他作为一封信的携带者去见你的善心。请友善地接待他，并用\\WorkTitle{圣经}所称赞的款待尊敬他，不要忘记蔬菜。\par

% structure: p0097 source=h3 target=subsubsection reason=relative-heading-rank; base=h2

\subsubsection{《书信》第六十二封}

\\EditorialNote{（这里提到的亚美尼亚人很可能是指\\Place{Sebaste}的主教\\Person{Eustathius}，该城是\\Place{小亚美尼亚}的首府。他曾是\\Person{Arius}的门徒，但多次公开宣认尼西亚信经，随其同伴改变意见。然而，他的个人品格非常高，很长一段时间，圣\\Person{Basil}以深情敬意看待他。事实上，圣\\Person{Basil}的隐修士规则正是基于他起草的一份规则。但在\\Person{Basil}升任主教后，\\Person{Eustathius}开始反对他，并在各方面诽谤他，甚至公开与亚略派共融。看来此人试图拉拢\\Person{Amphilochius}到他一边，并通过他拉拢\\Person{Gregory}。）}\par

你无与伦比的尊敬的训令不是野蛮的，而是希腊的，更确切地说是基督徒的；但你所引以为傲的那亚美尼亚人，却是一个彻头彻尾的野蛮人，远离我们的尊敬。\par

% structure: p0100 source=h3 target=subsubsection reason=relative-heading-rank; base=h2

\subsubsection{《书信》第六十三封。致\\Person{Amphilochius}（老）}

\\EditorialNote{（公元374年，\\Person{Amphilochius}被任命为\\Place{Iconium}的主教；他的父亲，同名的人，因被剥夺了儿子而深感悲痛，他曾指望儿子在他年老时供养他，并指责\\Person{Gregory}是原因。\\Person{Gregory}刚刚失去自己的父亲，写信来纠正他的错误认识，并使他相信，他多么害怕朋友和自己承担主教职责的重担。）}\par

您悲伤吗？我，当然，充满喜乐！您哭泣吗？我，如您所见，正在庆祝，并因目前的事态而夸耀！您是否因您的儿子被取走，并因他的德行而被提升尊荣而悲伤？您是否认为他不再与您在一起，照顾您的晚年，并按他的习惯给予您一切应有的关心和服务，是巨大的不幸？但对我来说，我的父亲踏上了最后的旅程，不再回到我这里，我再也不能见到他，这并不令我悲伤！那么，我这边并不责备您，也不要求您应有的慰问，因为我知道，私人的痛苦不容许为陌生人留出闲暇；没有人如此友善和如此哲学化，能超越自身的痛苦，并在自己需要安慰时安慰他人。但您反而雪上加霜，当我听说您责备我时，您认为您的儿子，也是我的兄弟，被我们忽视，甚至是被我们出卖，这是更重的指控；或者我们不承认他所有的朋友和亲属所受的损失，而我比所有人都更甚，因为我将生活的希望寄托在他身上，并将他视为唯一的堡垒、唯一的良善顾问、唯一与我共享虔诚的人。然而，您凭什么形成这种意见？如果是前者，请放心，我特意来到您这里，因为我被谣言困扰，并且我准备在事态还能商议时与您共同商议；但您却向我透露了别的，而不是这件事，不知是因为您也处于同样的痛苦中，还是有其他目的，我不知道。但如果是后者，我因为自己的悲伤、我对父亲应尽的荣誉以及他的葬礼（我不能让任何事超越它）而未能再次与您会面，而且当时我的悲痛刚刚发生，不合时宜地哲学化、超乎人性不仅是错误的，也是极不恰当的。此外，我认为我先前的处境已经够我忙碌的了，尤其是他的结局已按照那管理我们万事者的旨意成就。关于此事就谈这么多。现在我恳求您放下您的悲伤，我确信这极其不合理；如果您有任何进一步的委屈，请提出来，免得您部分地使我悲伤，也使您自己悲伤，并使您处于与您高贵身份不符的境地，责备我而非他人，而我并未对您行不义，但（如果必须说实话）被我们共同的朋友同样地压迫，尽管您曾认为我是您唯一的恩主。\par

% structure: p0103 source=h3 target=subsubsection reason=relative-heading-rank; base=h2

\subsubsection{Ep. CLXXI. 致\Place{依科尼雍}主教\Person{安斐洛邱}。}

我刚刚从病痛中解脱出来，便赶向您，我痊愈的守护者。因为一位思慕主的司铎的舌头能扶起病人。那么，在您司铎的服务中做更伟大的事，当您手持复活之祭时，解除我众多罪恶的重担。因您的事，无论醒着还是睡着，都是我的挂虑；您在我心中是一把好的拨子，在我灵魂中安置了一架和谐的里尔琴，因为您通过许多书信训练了我的灵魂走向科学。但是，最可敬的朋友，当您用您的言语牵引圣言降下，当您用血不流的切割分离主的身体和血，以您的声音作刀剑时，请不要停止为我祈祷和求情。\par

% structure: p0105 source=h3 target=subsubsection reason=relative-heading-rank; base=h2

\subsubsection{Ep. CLXXXIV.}

\EditorialNote{波索里乌斯，第二卡帕多西亚科洛尼亚的主教，他显然在选举和祝圣欧拉里乌斯为纳齐盎教座时起了重要作用，后被凯撒勒雅总主教赫拉迪乌斯指控为异端，公元383年在帕纳塞斯召开了一次会议审理此案。额我略因故不能亲自参加这次主教会议，写信给安斐洛邱，请求他担任被告的辩护人。该信已遗失，但额我略的朋友成功完成了他的使命，以下这封信是为了感谢他的好意。}\par

愿主满足您的一切祈求（不要轻视一位父亲的祈祷），因为您极大地安慰了我的晚年，既因为您如受邀前往了帕纳塞斯，又因为您驳斥了针对那位至可敬且天主所爱的主教的诽谤。因为恶人喜欢将自己的过失归于谴责他们的人。因这位教长的年纪比所有指控都更有力，他的人生也是如此，我们也常从他那里听到并教导他人，还有那些被他从错误中挽回并加入教会共同身体的人；但目前的邪恶时节要求更准确的证明，因为诽谤者和恶意的存在；而您为我们提供了这些，或者更确切地说，您为那些心意摇摆、容易被这样的人迷惑的人提供了这些。但如果您愿意走更远的路，并亲自作证，与其他主教一起解决此事，您将做一项配得上您成全的属灵工作。我和与我在一起的人问候您的兄弟情谊。\par

% structure: p0108 source=h2 target=subsection reason=relative-heading-rank; base=h2

\subsection{六、致\Place{君士坦丁堡}总主教\Person{内克塔留斯}}

\EditorialNote{额我略未能成功说服公元381年的大公会议通过承认保利诺为默勒提乌斯的继承人来结束安提约基雅的分裂，因此他认为为了和平最好辞去总主教职务。大公会议选举了内克塔留斯接替他，当时内克塔留斯是一位慕道者，并担任君士坦丁堡总督，他于381年6月9日被祝圣并登基。额我略始终与他保持友好关系；以下这封信是对他当选的正式通知的回应。}\par

% structure: p0110 source=h3 target=subsubsection reason=relative-heading-rank; base=h2

\subsubsection{Ep. LXXXVIII.}

王城理应有王像为其装饰。为此，她将你佩戴于胸前，这是理所当然的，连同你的德行、辩才，以及天主恩宠所显著赋予你的其他美善。但她却视我如敝屣，将我如同渣滓、糠秕或海浪一般抛弃。然而，既然朋友彼此休戚相关，我自当分享你的福祉，且自觉有份于你的光荣及其他一切顺遂。同样，你也应理所当然地分担你这流亡者的焦虑与挫折，不仅（如悲剧诗人所言）固守安乐，更要在患难中与你的朋友分忧；如此，你才能成为完全正直的人，在友谊中并对待朋友时生活得公正、平等。愿好运长伴你左右，使你行更多的善事；是的，愿它与你同在，永不更改，直到永远，在你今世的顺遂之后，直至进入那另一个世界。\par

% structure: p0112 source=h3 target=subsubsection reason=relative-heading-rank; base=h2

\subsubsection{Ep. XCI.}

\EditorialNote{此信无甚重要，唯可见\Person{国瑞}对其继任者持续保持友好之情。}\par

我们这边的情形照旧：我们安宁静谧，无争无辩，珍视沉默的酬报（此事毫无风险）胜过一切。我们亦从这安息中有所获益，因天主的仁慈，我们的病体已颇有好转。愿你策马前行，执掌王权，正如圣\Person{达味}所言；愿那以司祭职尊荣你的天主，始终与你同行，并为你高举它于一切毁谤之上。为使我们彼此能证明勇气，并在天主面前不受任何人世的灾厄，我以此信致你，愿你欣然应允。有许多理由使我极其挂念我们至为亲爱的\Person{潘克拉提乌斯}。请惠予善待他，将他推荐给你最好的朋友，使其得以达成目标。他的目标是藉某种军役摆脱公职，尽管没有任何一种生活方式能免于恶人的诽谤，这一点你深知。\par

% structure: p0115 source=h3 target=subsubsection reason=relative-heading-rank; base=h2

\subsubsection{Ep. CLI.}

\EditorialNote{写于约公元382年，为将其友、\Place{纳齐盎}的执事\Person{乔治}，因其私人的困扰与忧虑，推荐给总主教及宫廷总管或皇家内务总管之善职。}\par

众人对阁下的秉性大体上估计甚佳——更确切地说，他们并非猜测，而是当我自诩您认为我值得相当的尊重与荣誉时，他们并不拒绝相信。其中有我至为宝贵的儿子\Person{乔治}，他遭受诸多损失，为烦恼所困，惟见一个安全之港，即藉由我们引介到您面前，并从最尊贵的宫廷总管大人处获得一些恩惠。请将此恩惠赐予，或为他及其所需，或若你更愿如此，就当是为我，因我知道你已决心对我施以一切恩惠；事实也使我确信这一点。\par

% structure: p0118 source=h3 target=subsubsection reason=relative-heading-rank; base=h2

\subsubsection{Ep. CLXXXV.}

\EditorialNote{参见上文第四百六十九页Ep. CLXXXIV之序言。\Person{博斯波里乌斯}将被送往\Place{君士坦丁堡}，以便其案件在民事法庭审理。因此，\Person{国瑞}致信总主教，指出此举将严重侵犯教会的权利。很可能，\Person{涅克塔里乌斯}在\Person{国瑞}建议下所采取的态度，促成了\Person{德奥多西}于公元384或385年二月致奥古斯塔利斯长官的谕令，将所有神职人员从民事法庭的管辖中撤出，置于主教法庭的专属控制之下。}\par

每当不同的人称赞你的不同优点，并如市集般争先传扬你的美名时，我也尽己所能，且不逊于任何人，因为你也屈尊尊荣我，如爱子安慰其父一般安慰我的暮年。为此，我现在也斗胆向你呈上此恳求，为至为可敬、蒙主眷爱的主教\Person{博斯波里乌斯}；一方面，我羞于这样的一个人竟需要我的信函，因他令人起敬的品格已由其日常生活及年岁所证实；另一方面，我亦同样羞于沉默不语，不为他说一句话——我既有言说之能，又珍视信德，且与此人至为熟稔。关于教区之争，你无疑将凭你内住的圣神之恩宠及教规之秩序自行解决。但我愿阁下明鉴，我们的事务绝不可被张贴于世俗法庭。因为即使审理此类案件的法官是基督徒（如天主仁慈所赐），刀剑与圣神有何相干？即使我们让步，关乎信仰的争议又怎能、何处可以被公正地与其他问题交织？难道我们蒙主眷爱的主教\Person{博斯波里乌斯}今日成了异端？难道今日他那已使如此多人脱离谬误、已给出如此正统信仰之明证、并为我们众人之导师的皓首，却置于天平之上？不，我恳求你，切勿给这样的毁谤留地步；若可能，请使对立的双方和解，将此增添于你的赞美；若此不可行，至少不要容许我们众人（与他一同生活、一同老去之人）被这样的狂妄所凌辱——你知道我们是福音准确的宣讲者，无论是在宣讲危险之时，还是安全无虞之时；并且我们无法忍受任何对唯一不可接近之神性的诋毁。又求你为我这身患重病的人祈祷。我和我身边的所有人都问候围绕你的弟兄们。愿你强壮、勇敢，在主内享有美名，并将那众人共同要求的支持赐予我和教会。\par

% structure: p0121 source=h3 target=subsubsection reason=relative-heading-rank; base=h2

\subsubsection{Ep. CLXXXVI.}

\EditorialNote{一封为亲戚写的介绍信。}\par

如果我亲自前来占用您的时间，您会怎么做？我十分确定，您会以全部的热忱承担起为我洗刷诽谤的责任，如果我可以将先前发生的事视为一个征兆的话。那么，就通过我那位最谨慎的女亲属——她通过我来接近您——为我施此恩惠吧：首先，尊敬您这位求告者的年龄；其次，尊敬她的品德和虔诚，这超乎寻常地见于女性；除此之外，还有她对事务的无知，以及现在她的亲属给她带来的烦恼；尤其是我的恳求。您能给我的最大恩惠，是我所请求的那件事的迅速办理。因为福音中那个不义的判官，虽然是在长久恳求和纠缠之后，才向寡妇显示了仁慈。但我向您请求的是迅速，以免她因长期在异国他乡被焦虑和痛苦所压倒；尽管我很清楚，您的虔敬将使那异乡成为她的祖国。\par

% structure: p0124 source=h3 target=subsubsection reason=relative-heading-rank; base=h2

\subsubsection{书信第二百零二号}

（关于阿波利纳争议的一封重要信件已在上文给出。）\par

% structure: p0126 source=h2 target=subsection reason=relative-heading-rank; base=h2

\subsection{七、致\Place{提雅纳}主教\Person{德奥多禄}}

（\Person{德奥多禄}，\Place{阿利安祖斯}本地人，\Person{额我略}的亲密朋友，于公元379年陪同他前往\Place{君士坦丁堡}，并分担了亚略派人士的迫害，他们在举行神圣礼仪时闯入他们的教堂，用石头砸向神职人员。\Person{德奥多禄}无法容忍此事，并宣布打算起诉肇事者。\Person{额我略}写了以下信件劝阻他这样做，向他表明宽恕比报复更高贵。）\par

% structure: p0128 source=h3 target=subsubsection reason=relative-heading-rank; base=h2

\subsubsection{书信第七十七号}

我听说你对隐修士和乞丐们对我们实施的暴行感到愤慨。这并不奇怪，因为你从未受过打击，没有经历过我们必须忍受的苦难，因此你对这样的事情感到愤怒。但我们经历过许多种类的邪恶，也受过侮辱，因此当我们劝告阁下时，或许值得相信，正如老年所教导和理性所建议的。当然，所发生的事是可怕的，甚至比可怕更甚——没有人会否认：我们的祭台受到侮辱，我们的奥迹受到干扰，我们自己在领圣体者和欲投石者之间站立，用我们的代祷来医治疗伤；对贞女的尊敬被遗忘，隐修士的良好秩序被遗忘，穷人的灾难被遗忘，他们因残暴甚至失去了怜悯。但也许最好忍耐，并通过我们的痛苦为许多人树立忍耐的榜样。因为论证对一般世界的说服力不如实践，即那无声的劝勉。\par

我们认为获得那些错待我们之人的惩罚是重要的事：重要，我说（因为这对纠正他人有时也有用）——但遭受伤害并承受之，则更加伟大和更肖似天主。因为前者抑制邪恶，但后者成就良善，这比仅仅不邪恶要好得多和完善得多。让我们思考仁慈的伟大追求摆在我们面前，让我们宽恕对我们所做的错事，使我们也能获得宽恕，让我们以仁慈储存仁慈。\par

\Person{丕乃哈斯}被称为\OriginalTerm{热诚者}{Zelotes}，因为他用枪刺死了与那行淫的男人一起的米德杨女人\BibleRef{户 24:7}，并因为他除去了以色列子民的耻辱；但他更受赞扬是因为他在百姓犯罪时为百姓祈祷。那么，让我们也站立并做赎罪，让瘟疫停止，让这算作我们的义德。\Person{梅瑟}也受赞扬因为他杀了压迫以色列人的埃及人\BibleRef{出 2:12}；但他更令人敬佩是因为他通过祈祷治愈了他的姐妹\Person{米黎盎}，她因抱怨而得了癞病\BibleRef{户 12:40}。再看看后面。\Place{尼尼微}百姓被威胁要倾覆，但他们用眼泪赎回了他们的罪\BibleRef{纳 3:10}。\Person{默纳舍}是最无法无天的君王\BibleRef{编下 33:12}，但他是最显著的通过哀伤获得救恩的人之一。\par

\BibleQuote{厄弗辣因，我应怎样待你呢？}\BibleRef{欧 6:4} \Person{天主}说。这里表现的是何等的愤怒——但保护也随之而来。还有什么比仁慈更迅速？门徒们要求从天上降下\Place{索多玛}的火焰，烧毁那些拒绝\Person{耶稣}的人，但\Person{耶稣}拒绝了报复\BibleRef{路 9:54}。\Person{伯多禄}割下了恶意对待\Person{耶稣}的\Person{马耳古}的耳朵，但\Person{耶稣}复原了它。还有，那问是否应该七次原谅得罪他的弟兄的人，难道不是因他的吝啬而被谴责吗？因为七次之外增加了七十个七次\BibleRef{玛 18:21}。福音中那个不原谅被宽恕的债户呢？难道不是因此而更严厉地追讨他吗？祈祷的模范怎么说呢？难道不是希望宽恕可以通过宽恕来赚得吗？\par

既有如此多的榜样，让我们效法天主的仁慈，不要从自身学习到报复罪过是何等大恶。你看到善行的序列：首先制定法律，然后发令、威吓、谴责、警告、约束、再次威吓，只有在不得已时才施加打击，但这打击是逐渐的，为改过自新开路。因此，我们不要突然击打（因为这并不安全），而是要在敬畏中自我克制，以仁慈战胜，以善待使他们成为我们的负债人，用他们的良心而非愤怒折磨他们。我们不要干枯一棵可能还会结果子的无花果树（\BibleRef{路 13:7}），也不要判定它为无用而白占土地，说不定一位有技巧园丁的照料和勤勉还能治愈它。我们不要因着或许来自魔鬼的恶意和仇恨，如此迅速地毁坏如此伟大光荣的工程；而应选择显明自己仁慈而非严厉，爱穷人而非抽象正义；不要更看重那些煽动我们这样做的人，胜过那些劝阻我们的人，至少应考虑与乞丐争辩的羞耻，他们有一个巨大的优势，即即使他们错了，也会因其不幸而获得怜悯。但如今，请考虑所有穷人以及支持他们的人，所有隐修士和贞女都伏在你脚下，为他们祈祷。请为他们赐予这一恩惠（因为他们已经受够了，从他们向我们请求之事可见），尤其赐予我，作为他们的代表。如果我们被他们羞辱，你或许觉得奇怪，但请记住，当我们向你提出这一请求时，若不被你倾听，那会更糟糕。愿天主宽恕高贵的\Person{保禄}对我们的侮辱。\par

% structure: p0134 source=h3 target=subsubsection reason=relative-heading-rank; base=h2

\subsubsection{第一一五书}

此书约于公元382年复活节期间寄出，附有一册\WorkTitle{斐罗卡里亚}，即由他自己与\Person{圣巴西略}编辑的\Person{奥力振}著作的文选。\par

您以热忱超前了节日、信件，甚至时间，赐予我们一个预庆。这就是阁下所赐予我们的。作为回报，我们给予您我们最伟大的东西，即我们的祈祷。但为使您有些小事物纪念我们，我们送给您一本\Person{奥力振}的\WorkTitle{斐罗卡里亚}，内中选录了对文学爱好者有用的段落。请惠予接受，并以勤勉和圣神的帮助，证明它的效用。\par

% structure: p0137 source=h3 target=subsubsection reason=relative-heading-rank; base=h2

\subsubsection{第一二一书}

此信稍后写成，为感谢一份复活节礼物。\Person{德奥多若}刚被任命为\Place{提亚纳}总主教。\par

我们喜乐于爱的标记，尤其在此季节，且来自一位既如此年轻又如此完美的人；并且，用圣经的话问候您：在青年时期坚定，因为圣经就是这样称呼那些比年龄所预期更富有智慧的人。古圣祖为他们的子女祈求天上的甘露和土地的肥沃（\BibleRef{创 27:28}）以及其他诸如此类的事物，虽然或许有人以更高的意义理解这些；但我们要以属灵的意义全部回报给您。愿主成就您所有的请求，并愿您成为您子女的父亲（若我可简洁而亲密地为祈求），正如您自己在父母面前所彰显的那样，好使我们以及众人因您而受到荣耀。\par

% structure: p0140 source=h3 target=subsubsection reason=relative-heading-rank; base=h2

\subsubsection{第一二二书}

您欠我照顾，如同我是一个病人，因为诫命中有一条是探望病人。您也欠圣殉道者年度的敬礼，我们于您自己的\Place{亚里安祖斯}在称为\OriginalTerm{达图撒月}{Dathusa}的第二十三日举行。同时还有不少教会事务需要我们共同检讨。因此种种，我恳请您立刻前来：因为虽然劳苦很大，但赏报也相称。\par

% structure: p0142 source=h3 target=subsubsection reason=relative-heading-rank; base=h2

\subsubsection{第一二三书}

为自己推迟接受邀请而道歉。\par

我尊敬您的莅临，喜悦与您为伴；尽管另一方面，我劝自己留在家中安静默想，因我发现这在一切途径中对我最为有益。既然风势仍有些猛烈，我的虚弱尚未离去，我恳请您耐心容忍我少许时日，并请与我一同为健康祈祷；一旦合适季节来到，我必遵命前来。\par

% structure: p0145 source=h3 target=subsubsection reason=relative-heading-rank; base=h2

\subsubsection{第一二四书}

稍后，天气较为稳定时，\Person{额我略}接受了邀请，并提议立即前往，但拒绝参加省区主教会议。\par

您召唤我？我即赶来，且是为私下拜访。主教会议和集会我从远处致意，因为我经验到它们大多数（说得客气些）只是令人遗憾的事。那还剩下什么？请以您的祈祷助我正当的愿望，使我得以获得我所渴望的。\par

% structure: p0148 source=h3 target=subsubsection reason=relative-heading-rank; base=h2

\subsubsection{第一五二书}

自他从\Place{君士坦丁堡}引退后，\Person{额我略}应省区主教们的请求，尤其是\Place{提亚纳}都主教\Person{德奥多若}和\Place{科洛尼亚}主教\Person{波司波里乌}（见上文信件），以及民众的热切恳求，承担了\Place{纳齐盎}教区的管理；但他很快就发现自己的健康不能胜任如此重大的任务，无法满足教区对他的要求。他坚持了一段时间，但最终发现自己实在难以担当，于是致函都主教如下：\par

现在是时候让我使用这些圣经的话了：\BibleQuote{我受冤屈时该向谁呼求？}（\BibleRef{哈 2:1}）我受压迫时谁会向我伸手？这教会在当前邪恶瘫痪的状态下，该由谁来承担它的重担？我在天主和蒙选的天使面前声明，天主的羊群正受到不公正的对待，因为缺少牧人或主教，而我却被打入冷宫。因为我被自己的健康不佳所困，并因此很快被排除在教会之外，对所有人都毫无用处，每天都在垂死呼吸，被职责压得越来越不堪。如果这教省有其他首脑，我就有责任不断向它呼喊和抗议。但由于尊驾是上司，我必须仰望您。因为，撇开其他一切不说，您将从我的同僚司铎、\OriginalTerm{乡村主教}{Chorepiscopus}厄瓦辣（\Person{Eulalius}）和塞莱乌斯（\Person{Celeusius}）那里得知——我特意派遣他们到尊驾那里——这些现已占上风的强盗正在做什么和威胁什么。压制他们非我软弱的力量所能，而是属于您的技巧和力量；因为天主在赐予您其他恩赐的同时，也赐予了力量来保护祂的教会。如果我在言说和书写中得不到倾听，我将采取我唯一剩下的办法，即公开宣告并周知，这教会需要一位主教，以免因我虚弱的健康而受损。接下来如何，有待您考虑。\par

% structure: p0151 source=h3 target=subsubsection reason=relative-heading-rank; base=h2

\subsubsection{第153封信：致\Place{科洛尼亚}主教\Person{波司玻里}}

（圣额我略不得不实施他的威胁。他辞去了\Place{纳西盎}的牧职，什么都不能使他撤回辞呈。\Person{波司玻里}为此给他写了一封紧急信件，但他回复如下：）\par

我两次被你绊倒，并受了欺骗（你知道我指的是什么）。如果是公正的，愿主从你那里闻到\BibleQuote{馨香之气}（\BibleRef{创 8:21}）；如果是不公正的，愿主宽恕它。因为我这样对你说既是合理的，因为命令我们要忍耐所受的伤害。但正如你对自己的意见是主人，我对自己的意见也是主人。那个麻烦的额我略将不再让你麻烦。我要退到天主那里，唯有祂是纯洁无邪的。我要隐退于自身之内。我已决定这样做；因为两次在同一块石头上绊倒，谚语说只有愚人才会如此。\par

% structure: p0154 source=h3 target=subsubsection reason=relative-heading-rank; base=h2

\subsubsection{第157封信：致\Place{提雅纳}总主教\Person{德奥多若}}

（圣额我略在382年年底成功地使都主教及其同省主教相信他退休的诚意；于是他们开始物色一位继任者。额我略希望提名他的堂弟，\OriginalTerm{乡村主教}{Chorepiscopus}厄瓦辣，但主教们认为他试图向他们发号施令，感到某种嫉妒，拒绝让他参与选举，理由是他要么从未成为，或至少已不再是该教省的主教之一。他提出抗议，但发现自己无法说服他们，便撤回了他的投票权，并写信给\Person{德奥多若}如下：——）\par

我们属灵的事已到了极限：我不会再麻烦你了。团结起来：采取预防措施：商议反对我们：让我们的敌人得胜：让教规被准确遵守，从我们这些最无知的人开始。准确中并无恶意；只是不要让友谊的权利受到阻碍。我非常尊贵的儿子\Person{尼科布鲁}的孩子们来到城里学习速记。请以慈父般和善意的眼光看待他们（因为教规并不禁止这点），但特别注意他们住在教会附近。因为我希望他们通过与阁下的持续交往，在品德上被塑造成德性。\par

% structure: p0157 source=h3 target=subsubsection reason=relative-heading-rank; base=h2

\subsubsection{第163封信}

（\Place{提雅纳}的\Person{德奥多若}派遣\Place{帕斯帕苏斯}的一位平信徒\Person{乔治}到圣额我略那里，以便后者能使他相信自己的错误和罪恶，因为他背弃了所起的誓言，其理由是该誓是以书面形式而非\Emph{口头}形式起的。额我略似乎使他恢复了更好的心智，并把他送回都主教处，附上以下信件，请求对他处以适当的补赎，补赎的时间长短由其痛悔程度决定。这封信在553年第二次君士坦丁堡大公会议上由\Place{提雅纳}主教\Person{德奥多若}的继任者\Person{厄乌弗兰特}宣读，并被教父们接受，因此它被认为几乎具有普世教会教规的效力。）\par

愿天主将您赐予众教会，既为我们的荣耀，也为众人的益处；您在灵性事务上如此谨慎周详，甚至使我们这些在年岁上自以为比您稍长的人也更加谨慎。然而，您既愿我们参与您的灵性探讨（我指的是关于帕斯帕苏斯的\Person{乔治}似乎曾立下的\OriginalTerm{誓愿}{oath}），我们将向尊座陈述我们心中的想法。许多人自欺，认为以口头咒诅为佐证的誓愿才是真正的誓愿，而书面写就、未曾口头念出的誓愿则不过是形式，根本不算誓愿。因为我们怎能设想，一张书面的债务清单比口头承认更具约束力，而书面的誓愿却不算誓愿呢？简而言之，我们认为誓愿是向请求并获得它的人所作的保证。仅仅说他受了强迫是不够的（因为他约束自己的法律就是那强迫），也不能说他后来在法庭上赢了官司——因为他去打官司这一事实本身就违反了他的誓愿。我已经说服了我们的弟兄\Person{乔治}，不要再为他的罪过寻找借口，也不要寻求论据来辩护他的过犯，而要承认那书面文件是一个誓愿，并在天主和尊座面前哀哭他的罪过，即使他先前曾自欺，持不同看法。这是我们本人与他辩论的内容；显然，您若与他更仔细地交谈，必会加深他的痛悔，因为您是伟大的灵魂医治者；您按\OriginalTerm{教规}{Canon}规定的时间处理他，待时机合宜后，便可在时间上给予他\OriginalTerm{大赦}{indulgence}。而时间的量度，将是他痛悔的量度。\par

% structure: p0160 source=h3 target=subsubsection reason=relative-heading-rank; base=h2

\subsubsection{第一八三书}

\Person{赫拉狄乌斯}，\Place{凯撒勒亚}总主教，质疑\Person{欧拉里乌斯}被选为\Place{纳齐盎}主教之职的有效性，并指控\Person{博斯波里乌斯}犯有异端。圣\Person{额我略}在此以其全部权威支持另一方。然而，从这封信的措辞可以明显看出，收信人并非\Place{提亚纳}的\Person{德奥多若}。\Person{克雷芒塞}推测，他或许是\Place{莫普苏埃斯提亚}的\Person{德奥多若}。\par

嫉妒，无人能轻易逃脱，已在我们中间有了立足之地。看，甚至我们\Place{卡帕多西亚}人也处于分党结派的状态——可以说——一种前所未闻且难以置信的灾祸，使凡有血气的，在天主面前，一个也不能自夸\BibleRef{格前 1:29}，好叫我们小心，既然我们都是人，不要彼此轻率定罪。就我而言，甚至从这不幸中也有所得（若可说得有些悖论），我真如谚语所说，从荆棘中摘取玫瑰。直到如今，我从未面见尊座，也未曾以书信与您交谈，仅因您的声誉而得光照；但现在我不得已要借书信接近您，我十分感激那赐我此殊荣的人。我略过不写信给您信中所提的其他主教，因为时机尚未成熟。而且我衰弱的身体使我在此事上较为消极；但我写给您的，也通过您写给他们。我主内可敬的主教\Person{赫拉狄乌斯}必须停止在我们的事务上浪费精力。因为他寻求此事并非出于灵性热心，而是出于党派热忱，并非为了精确遵守教规，而是为了满足怒气，如同他选择的时间所表明的，且因许多人无理性地与他一同行动——我必须说这话，并不为此烦恼。倘若我身体条件允许治理\Place{纳齐盎}教会——我原被任命于此，而非\Place{萨西玛}，有人错误地劝您相信——我不会如此怯懦，或如此不谙\OriginalTerm{神圣法令}{Divine Constitutions}，以致要么轻视那教会，要么寻求安逸生活，而不愿得到那些按天主旨意劳苦、善用交托给他们管理的塔冷通者所储存的奖赏。因为若我在最重要的事上失策，我从诸多劳苦和伟大希望中能得什么益处呢？但由于我身体健康不佳，众人有目共睹，而且我因上述理由不担心此次引退有任何责任，我见到教会因紧附于我而在最佳利益上受苦，几乎因我的疾病而毁灭，我曾祈求，且现在再次祈求我主内可敬的主教们（我指我们\OriginalTerm{教省}{province}的）为教会立一位首领，他们已藉天主的恩宠这样做了，符合我的愿望和你们的祈祷。我愿您，最可敬的主，既亲自知道这事，也告知其余主教，好使他们接纳他，并以投票支持他，不要因听信诽谤而加重我晚年的负担。我还要补充一点：若您审查发现我主内可敬的司铎\Person{博斯波里乌斯}在信仰上有罪——这是个甚至不该暗示的事——（我略过他的年龄和我个人的证言）你们自己就这样定他吧。但如果关于教区的讨论是这恶言和这新奇指控的原因，请不要被诽谤带走，也不要给予谎言比真理更大的力量，我恳求您，免得使那些愿行正义的人陷入绝望。愿您获得良好的健康、精神、勇气，在天主的事情上不断进步，为我们和教会，您是我们共同的夸耀。\par

% structure: p0163 source=h3 target=subsubsection reason=relative-heading-rank; base=h2

\subsubsection{第一三九书}

这封信写于稍早时期，内容涉及他被迫同意继续担任\Place{纳齐盎}教区管理者一段时间的事。该信可以肯定不是写给\Place{提亚纳}的\Person{德奥多若}的，而且也不知道这位\Person{德奥多若}是谁。\par

那一位曾将祂的仆人\Person{达味}从牧人的工作中提升到王座，又将尊座从羊群中提升到牧人的职务；祂按照自己的旨意安排我们和一切寄望于祂的人的事务：愿祂现在使尊座心中明了，我在诸位主教大人关于投票之事上所遭受的羞辱——他们同意了选举，却排除了我们。我不愿将责任归咎于尊座，因为您刚来主持我们的事务，按情理，大多不熟悉我们的历史。这就足够了：因为我不愿再劳烦您，以免在我们友谊之初就显得烦扰。但我要告诉您，我在与天主商议时心中所起的念头。我从\Place{纳齐盎}的教会退隐，并非因为轻视天主，或轻视小羊群（愿哲思的灵魂绝不如此）；首先是因为我不受任何此类任命的约束；其次是因为我因体弱多病而心力交瘁，认为自己不足以承担这些忧虑。既然您也因我的辞职而责备我，而我自身也无法忍受对我的非议，又因时势艰难，有敌人入侵的威胁，损害整个教会的公益，我终于下定决心忍受一次对我的身体来说是痛苦、但对灵魂或许并非不利的失败。我将这可怜的身体交给教会，尽可能长久地使用，认为宁可忍受肉身上的任何痛苦，也不愿使我自己或他人遭受灵性的伤害——那些人以自己的经历判断，对我们做了最坏的揣测。知道了这一切，请为我祈祷，并赞同我的决定：也许说这话并不不合时宜，要塑造自己趋向虔诚。\par

% structure: p0166 source=h2 target=subsection reason=relative-heading-rank; base=h2

\subsection{八、致\Person{尼科布鲁}}

（参见上述致\Person{索弗若尼乌斯}的第一封信的导言。）\par

% structure: p0168 source=h3 target=subsubsection reason=relative-heading-rank; base=h2

\subsubsection{书信第十二（约公元365年）。}

你拿\Person{阿利皮安娜}取笑，说她身材矮小，配不上你的高大——你这身材和力量都如此高大、魁梧、奇异的家伙！因为现在我明白，灵魂是可以用尺寸衡量的，德性是可以用重量称量的，岩石比珍珠更珍贵，乌鸦比夜莺更可敬。好极了！为你的高大和你的肘尺而欢喜吧，在任何方面都不逊于著名的\Person{阿洛欧斯}之子。你骑马、执矛、关心野兽。但她却不做这些事；无需多大的力量来拿梳子、执纺锤、坐在织机旁，\Quote{因为如此是妇女的光荣。}如果你再加上这一点：她因祈祷而扎根于地，以灵魂的伟大运动与天主不断共融，那么你的高大或身体的伟岸又有什么可夸耀的呢？注意适时的沉默：聆听她的声音：留意她的朴素，她的女性之勇，她在家庭中的有用，她对她丈夫的爱。那时你将和拉科尼亚人一同说：确实，灵魂不是尺寸的主题，外在的必须看向内在的人。你若这样看待这些事，就会停止取笑和嘲弄她的矮小，而因你的婚姻庆幸自己。\par

% structure: p0170 source=h3 target=subsubsection reason=relative-heading-rank; base=h2

\subsubsection{书信第五十一}

（答复\Person{尼科布鲁}请求一篇关于写信艺术的论文。\Person{贝努瓦}认为此信及以下信件是写给小\Person{尼科布鲁}的。）\par

在写信的人当中——既然你问这个问题——有些人写得过于冗长，另一些人则过于简短；两者都偏离了中道，好比射手瞄准靶子，有的射得太近，有的射得太远；因为偏离虽在相反的方向，却是同样的失误。书信的标准在于其有用性：既不能在无话可说时写得过长，也不能在话多时写得过短。什么？难道我们要用波斯尺或儿童的肘尺来衡量智慧，写得如此不完整，以至于不像在写，而像是在描摹正午的影子，或那些在你正前方交汇的线条——它们的长度被缩短，只以片段而非全貌呈现，只有通过某些端点才能辨认？我们在两方面都应避免过度，而取适中。关于简洁，这是我的意见；关于清晰，显然应尽可能避免演说风格，而更倾向于闲聊；简言之，最好的、最美的书信是能使无学问或有学问的读者都信服的：前者因其浅近，后者因其超越众人；而且它本身应当明白易懂。猜谜和诠释一封信同样令人不快。关于书信的第三个要点是优雅：如果我们不以任何枯燥乏味、或粗糙无修饰、或杂乱无章、或未经修剪（如他们所说）的方式写作，就能保证这一点；例如，风格缺乏格言、谚语、精辟之语，甚至笑话和谜语——这些能使语言甜美。然而，我们不能显得过度使用它们。完全省略它们显得粗俗，但滥用则显得贪得无厌。我们可以像在织物中使用紫色那样适度使用它们。修辞格我们将允许，但应少用且节制。对偶、平行句和工整的分句，我们留给诡辩家；如果我们有时也采用，那更多是戏谑而非认真。我最后的评论是一位聪明人关于鹰所说的：当鸟儿们选举国王，各自以各种装饰前来时，它最美的一点是它不认为自己美。这一点在写信中尤需注意：不靠外加的装饰，而尽可能自然。关于书信，我用一封信给你寄去这些；但或许你最好不要把这应用于我本人，因为我正忙于更重要的事。其余的你将自己领悟，因为你善于学习，而精通这些事的人会教你。\par

% structure: p0173 source=h3 target=subsubsection reason=relative-heading-rank; base=h2

\subsubsection{书信第五十二}

（\Person{尼科布鲁}请求\Person{格列高利}出版他的书信集。\Person{格列高利}寄去一份抄本。）\par

你是在向秋日草原求花，在涅斯托耳晚年为他武装，如今要求我写出辞藻华丽的文字，而我早已忽视了语言与生活中一切令人愉悦的事物。然而——既然你强加于我的并非欧律斯透斯或赫剌克勒斯式的苦役，而是一件十分惬意且安适的事情，即尽我所能为你收集我自己的书信——那么，请将这一卷置于你的书丛之中：一部并非情爱之作，而是演说之作，不是为了展示，更多的是为了实用，且是为了我们自家所用。因为不同的作者有不同的特点，或大或小。我的特点是，一有机会便以格言和正面陈述来教导。如同在亲生的孩子身上，在语言中也是如此，父亲总是清晰可见，正如同父母由身体特征所显现。我的特点正如我所提及的。你可以通过写作和从我所写的内容中获益来回报我。我不能要求或请求比这更好的报酬了，既对请求者更有益，也对给予者更相称。\par

% structure: p0176 source=h3 target=subsubsection reason=relative-heading-rank; base=h2

\subsubsection{第五十三书简}

（额我略将巴西略的书信集与自己的书信放在一起，并让它们占据首要位置。尼科布鲁斯对此似乎感到惊讶，并询问原因。额我略解释如下。）\par

我一直将伟大的巴西略置于我自己之上，尽管他持相反意见；如今我依然如此，这既是为了真理，也是为了友谊。这就是为什么我将他的书信置于首位，而将我的置于次位。因为我希望我们二人永远被联结在一起；同时，我也愿为他人提供谦逊与顺服的榜样。\par

% structure: p0179 source=h3 target=subsubsection reason=relative-heading-rank; base=h2

\subsubsection{第五十四书简}

论拉哥尼亚文风。简洁并非如你所想，只是写很少的字，而是在很少的字中说出很多内容。因此，我称\Person{荷马}极为简洁，而\Person{安提马科斯}冗长。何以故？因为我以内容而非字母数来衡量长度。\par

% structure: p0181 source=h3 target=subsubsection reason=relative-heading-rank; base=h2

\subsubsection{第五十五书简}

邀请。我追赶你时，你逃开；或许这是爱的法则，为使自己更有价值。那么，来吧，终于补满你长久拖延给我造成的损失。若家务事缠身，你日后可再离开我们，从而使自己成为更珍贵的渴慕对象。\par

% structure: p0183 source=h3 target=subsubsection reason=relative-heading-rank; base=h2

\subsubsection{第九封：致奥林匹乌斯}

（奥林匹乌斯于三八二年任第二卡帕多西亚总督。一封反对阿波利纳里派的信函已给出；其余将陆续呈上，主要是向他的庇护推荐不同的人。）\par

% structure: p0185 source=h3 target=subsubsection reason=relative-heading-rank; base=h2

\subsubsection{第一百零四书简}

我所领受的一切其他恩惠，我深知皆出于你的善意；愿天主以自己的仁慈为此报答你；愿其中一项就是，你能以善名和光辉从始至终履行你的总督职务。在我现在所请求的事上，我来此与其说是领受，不如说是给予，若如此说不算傲慢。我亲自将可怜的菲卢梅娜引荐给你，恳求你主持公道，并因她使我灵魂悲戚的泪水而感动。她本人会向你说明自已在何事上、被何人冤屈，因我不宜对任何人提出控告。但我有必要说这一点：寡妇和孤儿有权获得所有正直之人的援助，尤其是那些有妻子儿女的人——这些怜悯的伟大抵押，因为我们自己只是人，却受命审判人。请原谅我以信件为你代求，因我因健康状况不佳，无法亲见一位如此仁慈、如此德高望重的统治者，以至于你统治的开端都比他人任期结束时的美名更为珍贵。\par

% structure: p0187 source=h3 target=subsubsection reason=relative-heading-rank; base=h2

\subsubsection{第一百零五书简}

时光飞逝，斗争艰巨，我的疾病愈发沉重，几乎使我无法动弹。还剩下什么？只有向天主祈祷，并向你的善意恳求：一则，愿祂使你的心倾向更温和的谋略；二则，愿你不会生硬地拒绝我们的代祷，而是仁慈地接纳可怜的保禄——公道已将他置于你手中，或许是为使你能因你伟大的仁慈而更加显赫，并使我们（这样的）祷告在你面前蒙垂怜。\par

% structure: p0189 source=h3 target=subsubsection reason=relative-heading-rank; base=h2

\subsubsection{第一百零六书简}

这里有另一个人向你呈上一封信，若实话实说，你本人便是原因，因为你对他们的礼遇激发了这些信。这里又有另一位向你恳求的人，一个恐惧的囚徒，我们的同族尤斯特拉提乌斯，他同我们一起、并藉着我们恳求你的良善，因为他无法忍受永远反抗你的统治（即便合理的恐惧已令他心惊），也不愿经由除我以外的任何人恳求你，为的是使他能因使用这样的代祷者而更显明你对他的怜悯——你在任何情况下都因接受他们的恳求而使代祷者变得伟大。我将简说一事：其他一切恩惠你都赐予了我；但这恩惠你将赐予你自己的判断，因为你曾立意要以这样的荣誉安慰我的年迈与病弱。我再加上一句：你正不断使天主更加眷顾你。\par

% structure: p0191 source=h3 target=subsubsection reason=relative-heading-rank; base=h2

\subsubsection{第一百二十五书简}

（已见上文，§ 1。）\par

% structure: p0193 source=h3 target=subsubsection reason=relative-heading-rank; base=h2

\subsubsection{第一百二十六书简}

（当额我略在克桑塔里斯时，有机会见到奥林匹乌斯，但旧病复发使他未能利用此机会。他写信表达遗憾，并借此机会请求免除尼科布鲁斯管理驿站之责。）\par

我在梦中感到快乐。因为被带到隐修院附近以从沐浴中获得些许安慰，随后期望与您相遇，且这好运几乎已在掌握之中，却因耽搁了数日，突然被我的疾病击倒——这疾病在某些方面已经令人痛苦，在其他方面则预示着威胁。而且，若必须找出某种猜测来解释这不幸，我的遭遇与章鱼相似：它们若被强行从岩石上扯下，便会冒着失去附着于岩石的吸盘之险，或者带走岩石的一部分。我的情况也有几分如此。我原本若见到阁下就会向您请求的事，现在虽身处远方也敢于提出。我发现我的儿子\Person{尼科布鲁斯}因照料驿站以及密切留意隐修院而十分焦虑。他并非强壮之人，且极不喜独居。请随意支使他做其他任何事情，因为他渴望在一切事上为您的权威服务；但如果可能的话，请解除他的这一职责，即使没有其他理由，至少是为了尊崇他作为我的招待。既然我曾为许多人向您请求过许多恩惠，并且都得到了，现在我也需要您对我本人的仁慈。\par

% structure: p0196 source=h3 target=subsubsection reason=relative-heading-rank; base=h2

\subsubsection{第131封}

\EditorialNote{公元382年，额我略奉召前往君士坦丁堡参加主教会议；他写信给市政长官\Person{普罗科皮乌斯}，婉拒前往，理由是他极不喜爱主教会议，他说自己从未见过这些会议带来任何好处。然而，他似乎通过\Person{伊卡里乌斯}和\Person{奥林皮乌斯}收到了更紧急的召唤。他给\Person{伊卡里乌斯}的回信已遗失；给\Person{奥林皮乌斯}的回信如下。}\par

对我来说，比我的疾病更严重的是，没有人相信我病了，反而要我进行如此长途的旅行，将我推入我庆幸已退出的烦恼之中，我几乎以为应当为此感谢我身体的苦痛。因为安静与摆脱事务比忙碌生活的辉煌更为宝贵。我昨日已将此事写信给最显赫的\Person{伊卡里乌斯}，从他那里我收到了同样的召唤；现在我也恳请您的宽宏大量为我写这封信，因为您是我健康状况非常可靠的见证。我无能的另一个明证是，我甚至无法前来享受与您的交往，您是一位如此仁慈的总督，在美德方面如此卓越，以至于您任期之始的序曲都比别人任期结束时获得的好名声更光荣。\par

% structure: p0199 source=h3 target=subsubsection reason=relative-heading-rank; base=h2

\subsubsection{第140封}

我又写信，本该亲自来的；但我是从您本人那里获得这样做的勇气，啊，精神事务的仲裁者（先谈最重要的），以及公共福祉的矫正者——两者皆出于天主的眷顾；您也因您的虔诚而获得回报，使您的事务按您的心意兴盛，并唯有您能做到他人所无法企及的事。因为智慧与勇敢引领您的治理，前者发现该做的事，后者轻松执行所发现的事。而最伟大的是您用以指导一切的双手的纯洁。您的不义之财在哪里？从未有过；您最先将其流放，如同一个无形的暴君。恶意在哪里？已被定罪。偏袒在哪里？在此您确实有些让步（因为我要稍加责备您），但这是在效仿天主的仁慈，目前您的士兵\Person{奥雷利乌斯}通过我向您恳求这种仁慈。我称他为愚蠢的逃亡者，因为他把自己交在我们手中，并通过我们交在您手中，以我们的白发和司祭职（您曾多次表示对此的尊敬）作为庇护，如同躲在某个皇帝圣像之下。看，这只献祭的、不沾染鲜血的手将他带到您面前；这只手曾多次写下赞美您的话，我相信若天主延长您的任期——我说的是您和您的同僚\Person{特弥斯}——还会写更多。\par

% structure: p0201 source=h3 target=subsubsection reason=relative-heading-rank; base=h2

\subsubsection{第141封}

\EditorialNote{纳齐盎的居民在某种程度上丧失了公民权利；\Person{奥林皮乌斯}签署了剥夺该城称号的命令。这导致部分年轻公民几乎叛乱；\Person{奥林皮乌斯}决定以彻底摧毁该地作为惩罚。圣额我略再次因生病而无法亲自出现在总督面前；但他通过以下信件为其家乡城（使用其正式拉丁名\Place{凯撒利亚}）辩护，并成功说服\Person{奥林皮乌斯}赦免了这次暴动。}\par

又一次提供仁慈的机会：我又大胆以书信托付我对如此重要之事的恳求。我的疾病使我如此大胆，因为它甚至不允许我外出，也不允许我向您做恰当的觐见。那么，我的使节任务是什么？请温和而仁慈地接受它。一个人的死亡固然可怕，他今日存在，明日便不复存在，不再回到我们中间；但一座城的死亡更为可怕，这座城是君王所建，时间所凝聚，漫长的岁月所保存。我说的是\Place{凯撒利亚}，曾经是一座城，如今不再有城，除非您赐予仁慈。请想，这个地方如今通过我俯伏在您脚下；让它拥有声音，披上丧服，如同悲剧中那样剪去头发，并用这样的话对您说：\par

\Quote{向我这个躺在尘埃中的伸出援手：帮助这个软弱无力者：不要将您的手的重量加在时间之上，也不要摧毁波斯人留给我的东西。对您而言，兴建城市比摧毁困境中的城市更光荣。请您做我的建立者，要么增加我所有，要么保持我原样。不要容忍直到您的治理时期我是一座城，而您之后便不再是城；不要给后世留下诽谤您的口实，说您接受我时我被列于城市之中，而您离开时我成了一片无人居住的荒地，曾经是城，如今只凭山脉、悬崖和树林才能认出。}\par

让我想象中的城市对您的仁慈如此说话和行事。但请屈尊接受我作为朋友的劝诫：您当然应惩罚那些反抗您权柄谕旨的人。在这一点上，我不敢多言，尽管据说这种鲁莽并非普遍策划，而只是少数年轻人的不理智愤怒。但请息怒大半，运用更宽广的理性。他们为母亲的被杀而悲伤；他们无法忍受被称为公民却没有政治权利；他们疯狂了；他们触犯了法律；他们抛弃了自己的安全；灾祸的突然使他们失去理智。难道为此城市就真的必须不再成为城市吗？当然不是。最卓越者，请不要为这件事写下命令。相反，请尊重所有公民、政治家及有身份之人的恳求——因为请记住，灾难将同样触及所有人——即使您伟大的权柄使他们保持沉默，仿佛暗中叹息。也请尊重我的白发：因为对我来说，拥有过一座伟大的城市，如今却一无所有，而且经过您的统治，我们为天主兴建的圣殿以及我们对装饰它的热爱，竟要成为野兽的居所，这将是可怕的。如果一些雕像被推倒，那并不可怕——尽管就其本身而言是可怖的——但我不希望您以为我在谈论这个，因为我全部关切在于更重要的事：但可怕的是，一座古城要与它们一同毁灭——它曾辉煌地存续，而我，受您尊敬且被认为略有影响力的人，竟活着看到了这一天。但关于此事就此足够，因为若我说得更多，也找不到比您自己的理性更强有力的东西——这个民族由您的理性治理，但愿更多更大的民族也受其治理，并在更大的统辖中。然而，有必要让阁下了解那些匍匐在您脚下的人，他们全然悲惨绝望，并未与那些违法者同流合污，正如许多当时在场者向我证实的那样。因此，请思量您认为适宜的决定，既为您在现世的名誉，也为您对来世的盼望。我们将承受您的决定——诚然不无悲伤——但我们会承受：因为我们还能做什么？若坏的决定占上风，我们将愤慨，并为我们已不复存在的城市洒落眼泪。\par

% structure: p0206 source=h3 target=subsubsection reason=relative-heading-rank; base=h2

\subsubsection{第一四二封信}

虽然我渴望见您的心意炽热，您的请愿者需求也迫切，但我的疾病无法克服。因此我冒昧将代祷托付于文字。请尊重我的白发，您已多次以善行予以尊敬。也请尊重我的软弱，我为天主劳苦在某种程度上促成了它——请容我稍加夸耀。为此，请宽恕那些仰望我的公民，因为我敢向您直言。也请宽恕我照管的其他众人。因为公共事务不会因仁慈受损，您藉畏惧所能成就的，胜过他人藉惩罚。愿您，作为赏报，得遇这样一位审判者，正如您向您的请愿者及我这代祷者所显示的那样。\par

% structure: p0208 source=h3 target=subsubsection reason=relative-heading-rank; base=h2

\subsubsection{第一四三封信}

丰富的经验，尤其是善的经验，为人带来什么？它教导仁慈，使人倾向于恳求者。没有什么教育比先前领受的良善更能教导怜悯。我自己就曾经历此事。我因所遭受的学会了同情。您看见我的心胸何等宽广：我为自己之事需要您的温和，却为他人代祷，并不怕耗尽您对他人事务的全部仁慈。我如此写信是为\Person{莱昂提乌斯}司铎——或者，若允许我如此称呼他，前任司铎。若他已为自己所行之事受罚足够，就让我们止步于此，免得过分成为不公。若仍有任何惩罚的余额，他罪行的后果尚未与他的冒犯相抵，但仍请为我们、为天主、为圣所，以及为司铎全体会议的缘故而赦免它——他曾被列于其中，尽管如今他因自己所行与所受，已显出自己不配。若我能说服您，那最好；若不能，我将带一位更有力的代祷者到您面前，她正是您统治与美名的伴侣。\par

% structure: p0210 source=h3 target=subsubsection reason=relative-heading-rank; base=h2

\subsubsection{第一四四封信}

（\Place{纳齐盎}公民\Person{维里亚努斯}被女婿冒犯，因此希望其女提起离婚诉讼。\Person{奥林匹乌斯}将此案交由\Person{圣额我略}主教仲裁，\Person{圣额我略}拒绝支持这一诉讼，并写下以下两封信：第一封致总督，第二封致\Person{维里亚努斯}本人。）\par

急躁并非总值得称赞。为此，关于最可敬的\Person{维里亚努斯}之女的事，我直到现在才作答复，既为给时间纠正事情，也因我猜想阁下不赞同离婚，因您将调查委托给我——您知道我在此类事上既不仓促也不冒失。因此我克制至今，我想并非没有理由。但既然我们已接近指定期限的尽头，您必须得知审查结果，我将告知您。在我看来，这位年轻女士似乎心有两念，介于对父母的尊敬与对丈夫的爱之间。言语站在他们一边，但我更认为心意与丈夫同在，她的眼泪便证明了这一点。您将按您正义所感到的，以及按在一切事上引导您的天主的旨意而行。我本最乐意向我儿\Person{维里亚努斯}进言，要他忽略许多争议之处，以避免确认离婚——这完全违反我们的法律，尽管罗马法可能另有规定。因为必须遵守正义——我祈祷您永远宣讲并实践正义。\par

% structure: p0213 source=h3 target=subsubsection reason=relative-heading-rank; base=h2

\subsubsection{第一四五封信。致\Person{维里亚努斯}。}

公共行刑者不犯罪，因为他们是法律的仆人；我们用以惩罚罪犯的剑也并非不合法。然而，公共行刑者并非可赞扬的角色，承载死亡的剑也非被欣然接受。同样，我无法忍受因用自己的手和言语确认离婚而招致憎恨。促成联合与友谊远比导致分裂与生命分离更好。我想，我们可敬的总督正是心中存有这一点，才将关于你女儿的调查托付给我，认为我不会仓促或冷酷地走向离婚。因为他提议我不是作为法官，而是作为主教，并将我置于你不幸境况中的调解人。因此，我恳请你多少包容我的胆怯；如果更好的方面占上风，就使用我作为你愿望的仆人：我乐于接受这样的命令。但若要采取更糟、更残忍的方式，就请找更适合你目的的人。我没有时间为了讨好你的友谊（尽管我在各方面都极其尊重你）而得罪天主，我必须向祂交账每一行为与思想。我相信你的女儿（因为真相终将道出），当她能够放下对你的畏惧，大胆宣告真相时。目前她的处境可悲——因为她将言语归于你，而将眼泪归于她的丈夫。\par

% structure: p0215 source=h3 target=subsubsection reason=relative-heading-rank; base=h2

\subsubsection{致\Person{奥林匹乌斯}。第一四六书。}

这就是我仿佛预言般所说的，当我发现你赞同每一个请求，并正不知餍足地使用你的温和时，我担心我会在其他人的事务上耗尽你的善意。因为看，我自己的争战临到了（如果那与我亲属有关的事算是我的），而我不能以同样的自由说话。首先，因为这是我自己的事。因为为自己恳求，尽管可能更有益，却更令人羞愧。其次，我担心过度会破坏愉悦，并反对一切美好。情况如此，而我的猜测只能说是太过正确。尽管如此，我信赖天主——我站在祂面前——并信赖你在行善上的宽宏大量，我大胆呈上此请愿。\par

假设\Person{尼科布卢斯}是最坏的人：——尽管他唯一的罪过是，因着我，他成为嫉妒的对象，且比应有的更自由。并假设我当前的对手是最正义的人。因为我不愿意在您正直的胸怀前控告一个我昨天还在支持的人：但我不知道，是否在您看来，因一个人的罪行而从另一个人身上索求惩罚是正义的，尽管这些罪行与后者完全无关，甚至未经他的同意；从那个搅动自己家宅、如此不安以至于比控告者所愿更轻易地投降的人身上。难道\Person{尼科布卢斯}或其子女必须如其迫害者所愿被贬为奴隶吗？我为迫害的理由和时机都感到羞愧，如果这事要在您掌权且我对您有影响力时发生。切勿如此，最可敬的朋友，不要让这事被呈到您的正直面前。但要藉着您思想的飞速敏捷，认出此事所出自的恶意，并顾及我——您的仰慕者，向那些被扰乱的人显示您仁慈的法官——因为今天您不仅是在人与人之间审判，而是在美德与恶习之间；对这种事，像您这样有德且精于治理的人必须比常人给予更多考虑。作为回报，您从我这里将不仅得到我的祈祷——我知道您不像许多人那样轻视它——而且我将在所有认识我的人中使您的治理闻名。\par

% structure: p0218 source=h3 target=subsubsection reason=relative-heading-rank; base=h2

\subsubsection{第一五四书。}

对我来说，即使您的任期已满，您仍是总督——因为我判断事物与常人不同——因为您在自己身上包含了一切总督的美德。因为许多坐在高耸宝座上的人，在我看来是卑下的，所有那些因自己的手而卑下并成为臣民奴隶的人。但许多人是崇高而高贵的，即使他们身处低位，美德将他们置于高处，并使他们配得上更大的治理。但这与我何干？伟大的\Person{奥林匹乌斯}不再与我们同在，也不再执掌我们的舵绳。我们完了，我们被出卖了，我们再次成为\Place{第二卡帕多细亚}，而您曾使我们成为\Place{第一卡帕多细亚}。其他人的事我为何要说？但谁来珍惜您的\Person{额我略}的晚年，用尊荣的魔力调治他的软弱，并因他从您那里为许多人谋得恩惠而使他更加尊贵？现在，请带着护卫和更大的排场启程吧，留下许多眼泪给我们，并带着大量财富——一种少数总督能拥有的财富：美名，以及被铭刻在众人心中，像不易移动的柱石。如果您再次以更大更光辉的统治治理我们（这是我们渴望的预兆），我们将向天主献上更完美的感谢。\par

\SourceRecord{Translated by Charles Gordon Browne and James Edward Swallow.；Nicene and Post-Nicene Fathers, Second Series；Vol. 7.；Edited by Philip Schaff and Henry Wace.；Buffalo, NY: Christian Literature Publishing Co.,；1894.；<http://www.newadvent.org/fathers/3103c.htm>.}
